\paragraph{零余类有限交截、nerve 的 flag 刚性、加权 clique--Euler 并集公式与共振整函数系数刚性}
在零点余类刚性、约数驱动的谱零点余类并结构、零点削减碰撞上界，以及共振梯整函数的 Laguerre--P\'olya 闭包、实零点结构与截断误差已经建立之后，还可继续沿同一条算术--解析主线抽取出下列七条结果。它们分别刻画零余类族任意有限交截的广义中国剩余判据、零点余类覆盖的 Helly--2 刚性与 nerve 的 flag 闭式、谱零点集合的精确加权 clique--Euler 公式、碰撞上界的约数兼容图闭式、共振整函数 \(\mathcal F_\varphi\) 的对数 Taylor 级数与倒数零矩、其偶系数的显式 Newton 递推与严格对数凹性，以及截断深度对相对误差预算的可审计控制。

\begin{theorem}[零余类族的任意有限交截由两两兼容性完全控制]\label{thm:xi-time-part71-zero-coset-finite-intersection-pairwise-compatibility}
\leanverified{paper\_xi\_time\_part71\_zero\_coset\_finite\_intersection\_pairwise\_compatibility}
设 \(M\) 为偶整数，且
\[
g_1,\dots,g_r\mid \frac{M}{2}.
\]
对每个 \(g\mid M/2\)，定义
\[
S_g(M):=
\left\{
k\in \ZZ/(M\ZZ):
\ k\equiv \frac{M}{2g}\pmod{M/g}
\right\}.
\]
则对任意有限族 \(\{g_1,\dots,g_r\}\)，都有
\[
\bigcap_{j=1}^{r}S_{g_j}(M)\neq\varnothing
\iff
2\gcd(g_i,g_j)\mid (g_j-g_i)
\qquad
(1\le i<j\le r).
\]
并且一旦交截非空，就有精确基数公式
\[
\left|
\bigcap_{j=1}^{r}S_{g_j}(M)
\right|
=
\gcd(g_1,\dots,g_r).
\]
\end{theorem}
\begin{proof}
由定理 \ref{thm:fold-zero-coset-rigidity}，
\[
S_{g_j}(M)
=
\left\{
k\in \ZZ/(M\ZZ):
\ k\equiv \frac{M}{2g_j}\pmod{M/g_j}
\right\}
\qquad (1\le j\le r).
\]
因此
\[
\bigcap_{j=1}^{r}S_{g_j}(M)\neq\varnothing
\]
当且仅当同余组
\[
k\equiv \frac{M}{2g_j}\pmod{M/g_j}
\qquad (j=1,\dots,r)
\]
可解。广义中国剩余定理断言：这一线性同余组可解，当且仅当任意两式相容。另一方面，引理 \ref{lem:fold-zero-coset-intersection} 已给出两式相容的充要条件恰为
\[
2\gcd(g_i,g_j)\mid (g_j-g_i).
\]
于是得到第一条等价关系。

现设该系统可解。由于解集模数为
\[
\operatorname{lcm}\!\left(\frac{M}{g_1},\dots,\frac{M}{g_r}\right)
=
\frac{M}{\gcd(g_1,\dots,g_r)},
\]
故在模 \(M\) 的意义下，解恰形成一个单一同余类的提升，其基数为
\[
\frac{M}{M/\gcd(g_1,\dots,g_r)}
=
\gcd(g_1,\dots,g_r).
\]
证毕。
\end{proof}

\begin{corollary}[零点余类覆盖的 Helly--2 刚性与 nerve 的 flag 闭式]\label{cor:xi-time-part71-zero-coset-nerve-flag}
沿用定理 \ref{thm:xi-time-part71-zero-coset-finite-intersection-pairwise-compatibility} 的记号。任取有限集合
\[
\mathcal G\subseteq
\left\{
g:\ g\mid \frac{M}{2}
\right\},
\]
并考虑覆盖
\[
\mathcal U_{\mathcal G}:=\{S_g(M):\ g\in\mathcal G\}.
\]
定义其 nerve 为
\[
\mathcal N(\mathcal U_{\mathcal G})
:=
\left\{
\sigma\subseteq \mathcal G:\ \bigcap_{g\in \sigma}S_g(M)\neq\varnothing
\right\}.
\]
再在顶点集 \(\mathcal G\) 上定义算术兼容图 \(\Gamma_{\mathcal G}\)：对不同的 \(g,h\in\mathcal G\)，规定
\[
g\sim h
\iff
2\gcd(g,h)\mid (h-g).
\]
则 \(\mathcal N(\mathcal U_{\mathcal G})\) 恰为 \(\Gamma_{\mathcal G}\) 的 clique 复形。特别地，\(\mathcal N(\mathcal U_{\mathcal G})\) 是一个 flag 复形：任意顶点集若两两相邻，则自动张成一个单纯形。
\leanverified{paper\_xi\_time\_part71\_zero\_coset\_nerve\_flag}
\end{corollary}
\begin{proof}
由定理 \ref{thm:xi-time-part71-zero-coset-finite-intersection-pairwise-compatibility}，对任意有限子集
\[
\sigma\subseteq \mathcal G
\]
都有
\[
\bigcap_{g\in \sigma}S_g(M)\neq\varnothing
\iff
2\gcd(g,h)\mid (h-g)
\qquad
(\forall\, g\neq h\in \sigma).
\]
右侧正表明 \(\sigma\) 在图 \(\Gamma_{\mathcal G}\) 中形成一个 clique。故
\[
\sigma\in \mathcal N(\mathcal U_{\mathcal G})
\iff
\sigma\in \mathrm{Cliq}(\Gamma_{\mathcal G}),
\]
即 nerve 恰为该兼容图的 clique 复形。于是它自动是 flag 复形。证毕。
\end{proof}

\endinput


\begin{theorem}[谱零点集合的精确加权 clique--Euler 公式]\label{thm:xi-time-part71-zero-spectrum-weighted-clique-euler}
\leanverified{paper\_xi\_time\_part71\_zero\_spectrum\_weighted\_clique\_euler}
设
\[
M:=F_{m+2},
\qquad
\mathcal Z_m:=\{k\in \ZZ/(M\ZZ):\ \widehat d_m(k)=0\},
\]
并假设 \(M\) 为偶数。定义
\[
\mathcal D_m:=
\left\{
d:\ d\mid (m+2),\ d<m+2,\ F_d\mid \frac{M}{2}
\right\}.
\]
在顶点集 \(\mathcal D_m\) 上定义兼容图 \(\Gamma_m\)：对不同顶点 \(d,e\in\mathcal D_m\)，规定
\[
d\sim e
\iff
2F_{\gcd(d,e)}\mid (F_e-F_d).
\]
若记 \(\mathrm{Cliq}(\Gamma_m)\) 为 \(\Gamma_m\) 的全部非空 clique，且对每个 clique \(C\) 记
\[
\gcd(C):=\gcd\{d:\ d\in C\},
\]
则谱零点集合满足严格公式
\[
|\mathcal Z_m|
=
\sum_{\varnothing\neq C\in \mathrm{Cliq}(\Gamma_m)}
(-1)^{|C|+1}F_{\gcd(C)}.
\]
\end{theorem}
\begin{proof}
由推论 \ref{cor:fold-zero-coset-union}，当 \(M\) 为偶数时，
\[
\mathcal Z_m
=
\bigcup_{\substack{1\le j\le m\\ g_j\mid M/2}}S_{g_j}(M),
\qquad
g_j=\gcd(F_{j+1},M)=F_{\gcd(j+1,m+2)}.
\]
去重之后，恰可写成
\[
\mathcal Z_m=\bigcup_{d\in\mathcal D_m}S_{F_d}(M).
\]

现取任意非空子集 \(I=\{d_1,\dots,d_r\}\subseteq \mathcal D_m\)。由定理 \ref{thm:xi-time-part71-zero-coset-finite-intersection-pairwise-compatibility}，
\[
\bigcap_{d\in I}S_{F_d}(M)\neq\varnothing
\]
当且仅当对所有 \(1\le i<j\le r\) 都有
\[
2\gcd(F_{d_i},F_{d_j})\mid (F_{d_j}-F_{d_i}).
\]
再由 Fibonacci 最大公因子恒等式
\[
\gcd(F_a,F_b)=F_{\gcd(a,b)},
\]
上式等价于
\[
2F_{\gcd(d_i,d_j)}\mid (F_{d_j}-F_{d_i}),
\]
也即 \(I\) 在 \(\Gamma_m\) 中形成一个 clique。

若交截非空，则再由定理 \ref{thm:xi-time-part71-zero-coset-finite-intersection-pairwise-compatibility}，
\[
\left|
\bigcap_{d\in I}S_{F_d}(M)
\right|
=
\gcd(F_{d_1},\dots,F_{d_r})
=
F_{\gcd(d_1,\dots,d_r)}
=
F_{\gcd(I)}.
\]
因此对并集
\[
\mathcal Z_m=\bigcup_{d\in\mathcal D_m}S_{F_d}(M)
\]
施行包含--排斥，便得到
\[
|\mathcal Z_m|
=
\sum_{\varnothing\neq C\in \mathrm{Cliq}(\Gamma_m)}
(-1)^{|C|+1}F_{\gcd(C)}.
\]
证毕。
\end{proof}

\begin{corollary}[碰撞上界可压缩为约数兼容图的加权 Euler 特征]\label{cor:xi-time-part71-collision-upper-bound-weighted-clique-euler}
沿用定理 \ref{thm:xi-time-part71-zero-spectrum-weighted-clique-euler} 的记号。则碰撞概率满足严格上界
\[
\mathsf{Col}_m
\le
1-\frac{1}{F_{m+2}}
\sum_{\varnothing\neq C\in \mathrm{Cliq}(\Gamma_m)}
(-1)^{|C|+1}F_{\gcd(C)}.
\]
特别地，若 \(\Gamma_m\) 不含三角形，亦即其非空 clique 仅有单点与边，则上界简化为
\[
\mathsf{Col}_m
\le
1-\frac{1}{F_{m+2}}
\left(
\sum_{d\in\mathcal D_m}F_d
-
\sum_{\{d,e\}\in E(\Gamma_m)}F_{\gcd(d,e)}
\right).
\]
\end{corollary}
\begin{proof}
由定理 \ref{thm:fold-collision-zero-reduction}，
\[
\mathsf{Col}_m\le 1-\frac{|\mathcal Z_m|}{F_{m+2}}.
\]
将定理 \ref{thm:xi-time-part71-zero-spectrum-weighted-clique-euler} 的精确公式代入，即得第一式。

若 \(\Gamma_m\) 不含三角形，则全部非空 clique 仅由单点与边组成，包含--排斥求和自动截断为两层：
\[
|\mathcal Z_m|
=
\sum_{d\in\mathcal D_m}F_d
-
\sum_{\{d,e\}\in E(\Gamma_m)}F_{\gcd(d,e)}.
\]
再代回碰撞上界即可。证毕。
\end{proof}

\begin{theorem}[共振整函数的对数 Taylor 级数与倒数零矩闭式]\label{thm:xi-time-part71-resonance-entire-log-taylor-zero-moments}
沿用定理 \ref{thm:fold-resonance-entire-lp} 的记号，记
\[
\mathcal F_\varphi(z):=\prod_{n=2}^{\infty}\cos(\pi z\varphi^{-n}).
\]
则在最近正零点半径
\[
|z|<\frac{\varphi^2}{2}
\]
内，有严格对数展开
\[
\log \mathcal F_\varphi(z)
=
-\sum_{r=1}^{\infty}\frac{\sigma_r}{r}\,z^{2r},
\qquad
\sigma_r:=
(2^{2r}-1)\zeta(2r)\frac{\varphi^{-2r}}{\varphi^{2r}-1}.
\]
等价地，正零点倒数偶矩满足
\[
\sum_{\substack{\rho\in\mathcal Z(\mathcal F_\varphi)\\ \rho>0}}\rho^{-2r}
=
(2^{2r}-1)\zeta(2r)\frac{\varphi^{-2r}}{\varphi^{2r}-1},
\]
而全零点求和式为
\[
\sum_{\rho\in\mathcal Z(\mathcal F_\varphi)}\rho^{-2r}
=
2(2^{2r}-1)\zeta(2r)\frac{\varphi^{-2r}}{\varphi^{2r}-1}.
\]
\end{theorem}
\begin{proof}
由定理 \ref{thm:fold-resonance-entire-lp}，\(\mathcal F_\varphi\) 为整函数，且
\[
\mathcal Z(\mathcal F_\varphi)
=
\left\{
\left(k+\tfrac12\right)\varphi^n:\ k\in\ZZ,\ n\ge 2
\right\}.
\]
当 \(|z|<\varphi^2/2\) 时，对所有 \(n\ge 2\) 都有
\[
|\pi z\varphi^{-n}|<\frac{\pi}{2},
\]
故可使用经典展开
\[
\log\cos x
=
-\sum_{r=1}^{\infty}\frac{(2^{2r}-1)\zeta(2r)}{r\pi^{2r}}\,x^{2r}
\qquad
(|x|<\pi/2).
\]
逐项代入 \(x=\pi z\varphi^{-n}\) 并交换求和，得到
\[
\log \mathcal F_\varphi(z)
=
\sum_{n=2}^{\infty}\log\cos(\pi z\varphi^{-n})
=
-\sum_{r=1}^{\infty}\frac{(2^{2r}-1)\zeta(2r)}{r}
\left(
\sum_{n=2}^{\infty}\varphi^{-2rn}
\right)z^{2r}.
\]
而
\[
\sum_{n=2}^{\infty}\varphi^{-2rn}
=
\frac{\varphi^{-4r}}{1-\varphi^{-2r}}
=
\frac{\varphi^{-2r}}{\varphi^{2r}-1},
\]
遂得
\[
\log \mathcal F_\varphi(z)
=
-\sum_{r=1}^{\infty}\frac{\sigma_r}{r}\,z^{2r}.
\]

另一方面，由定理 \ref{thm:fold-resonance-entire-lp} 的零点描述，
\[
\mathcal Z(\mathcal F_\varphi)\cap (0,\infty)
=
\left\{
\left(k+\tfrac12\right)\varphi^n:\ k\ge 0,\ n\ge 2
\right\}.
\]
故
\[
\sum_{\substack{\rho\in\mathcal Z(\mathcal F_\varphi)\\ \rho>0}}\rho^{-2r}
=
\sum_{n=2}^{\infty}\varphi^{-2rn}\sum_{k\ge 0}\left(k+\tfrac12\right)^{-2r}.
\]
又有
\[
\sum_{k\ge 0}\left(k+\tfrac12\right)^{-2r}
=
2^{2r}\sum_{k\ge 0}(2k+1)^{-2r}
=
(2^{2r}-1)\zeta(2r),
\]
故正零点公式成立。全零点关于原点成对对称，因此再乘以 \(2\) 即得最后一式。证毕。
\end{proof}

\begin{theorem}[共振整函数偶系数的显式 Newton 递推与严格对数凹性]\label{thm:xi-time-part71-resonance-entire-even-coefficient-newton-logconcavity}
\leanverified{paper\_xi\_time\_part71\_resonance\_entire\_even\_coefficient\_newton\_logconcavity}
将 \(\mathcal F_\varphi\) 写成偶幂级数
\[
\mathcal F_\varphi(z)=\sum_{n\ge 0}a_{2n}z^{2n},
\qquad
a_{2n+1}=0,
\]
并定义
\[
b_n:=(-1)^n a_{2n}
\qquad (n\ge 0).
\]
则对全部 \(n\ge 0\) 都有 \(b_n>0\)，并满足精确 Newton 递推
\[
n\,b_n
=
\sum_{r=1}^{n}(-1)^{r-1}\sigma_r\,b_{n-r}
\qquad (n\ge 1),
\]
其中 \(\sigma_r\) 如定理 \ref{thm:xi-time-part71-resonance-entire-log-taylor-zero-moments} 所示。进一步，序列 \((b_n)_{n\ge 0}\) 严格对数凹：
\[
b_n^2>b_{n-1}b_{n+1}
\qquad (n\ge 1).
\]
因此偶系数满足严格交替符号律
\[
a_0>0,\qquad a_2<0,\qquad a_4>0,\qquad \dots,
\]
且相邻比值 \((b_{n+1}/b_n)_{n\ge 0}\) 严格递减。
\end{theorem}
\begin{proof}
先由定理 \ref{thm:xi-time-part71-resonance-entire-log-taylor-zero-moments} 定义形式幂级数
\[
G(t):=\mathcal F_\varphi(\sqrt t)
=
\sum_{n\ge 0}a_{2n}t^n
=
\sum_{n\ge 0}(-1)^n b_n t^n.
\]
则
\[
\log G(t)
=
-\sum_{r=1}^{\infty}\frac{\sigma_r}{r}t^r.
\]
对 \(t\) 求导得
\[
\frac{G'(t)}{G(t)}
=
-\sum_{r=1}^{\infty}\sigma_r t^{r-1}.
\]
把
\[
G(t)=\sum_{n\ge 0}(-1)^n b_n t^n
\]
代入，并比较 \(t^{n-1}\) 的系数，即得
\[
n\,(-1)^n b_n
=
-\sum_{r=1}^{n}\sigma_r(-1)^{n-r}b_{n-r},
\]
等价于
\[
n\,b_n
=
\sum_{r=1}^{n}(-1)^{r-1}\sigma_r\,b_{n-r}.
\]

接着证明 \(b_n>0\) 与严格对数凹性。由定理 \ref{thm:fold-resonance-entire-lp}，
\(\mathcal F_\varphi\) 的正零点可记为
\[
0<\rho_1<\rho_2<\cdots,
\]
且由定理 \ref{thm:xi-time-part71-resonance-entire-log-taylor-zero-moments} 在 \(r=1\) 时的公式，
\[
\sum_{j\ge 1}\rho_j^{-2}<\infty.
\]
因此
\[
G(t)=\prod_{j=1}^{\infty}\left(1-\frac{t}{\rho_j^2}\right)
\]
在 \(t=0\) 邻域内成立，并且
\[
b_n=e_n(\rho_1^{-2},\rho_2^{-2},\dots)>0
\]
是正数列 \((\rho_j^{-2})_{j\ge 1}\) 的第 \(n\) 个初等对称函数。

对每个 \(N>n\)，定义截断多项式
\[
G_N(t):=\prod_{j=1}^{N}\left(1-\frac{t}{\rho_j^2}\right)
=
\sum_{s=0}^{N}(-1)^s b_s^{(N)}t^s.
\]
则 \(b_s^{(N)}\to b_s\)（\(s\) 固定）且 \(b_s^{(N)}>0\)。Newton 不等式给出
\[
\bigl(b_n^{(N)}\bigr)^2
\ge
\frac{n+1}{n}\frac{N-n+1}{N-n}\,
b_{n-1}^{(N)}b_{n+1}^{(N)}.
\]
令 \(N\to\infty\)，得到
\[
b_n^2
\ge
\frac{n+1}{n}\,b_{n-1}b_{n+1}
>
b_{n-1}b_{n+1}.
\]
这就证明了严格对数凹性。由于 \(b_n>0\)，于是
\[
a_{2n}=(-1)^n b_n
\]
严格交替。再将
\[
b_n^2>b_{n-1}b_{n+1}
\]
两边除以 \(b_{n-1}b_n>0\)，即可得到
\[
\frac{b_{n+1}}{b_n}<\frac{b_n}{b_{n-1}},
\]
故相邻比值严格递减。证毕。
\end{proof}

\begin{corollary}[共振梯截断深度的相对误差预算]\label{cor:xi-time-part71-resonance-truncation-relative-error-budget}
\leanverified{paper\_xi\_time\_part71\_resonance\_truncation\_relative\_error\_budget}
沿用命题 \ref{prop:fold-resonance-truncation-error} 的记号。对任意 \(\varepsilon>0\)，当 \(u=0\) 时有
\[
\frac{C_{\varphi,N}(0)}{C_\varphi(0)}=1
\qquad
(\forall N\ge 2).
\]
若 \(u\neq 0\)、\(N\ge 2\) 且
\[
N\ \ge\
-1+\max\!\left\{
\log_\varphi(2\pi|u|),\
\frac12\log_\varphi\!\left(
\frac{\pi^2u^2}{(1-\varphi^{-2})\log(1+\varepsilon)}
\right)
\right\},
\]
则有相对误差界
\[
1\le \frac{C_{\varphi,N}(u)}{C_\varphi(u)}\le 1+\varepsilon.
\]
\end{corollary}
\begin{proof}
当 \(u=0\) 时，
\[
C_{\varphi,N}(0)=C_\varphi(0)=1,
\]
结论平凡。以下设 \(u\neq 0\)。

由命题 \ref{prop:fold-resonance-truncation-error}，只要满足
\[
\pi|u|\varphi^{-(N+1)}\le \frac12,
\]
就有
\[
C_{\varphi,N}(u)\exp\!\left(
-\frac{\pi^2u^2\varphi^{-2(N+1)}}{1-\varphi^{-2}}
\right)
\le
C_\varphi(u)
\le
C_{\varphi,N}(u).
\]
因此
\[
1
\le
\frac{C_{\varphi,N}(u)}{C_\varphi(u)}
\le
\exp\!\left(
\frac{\pi^2u^2\varphi^{-2(N+1)}}{1-\varphi^{-2}}
\right).
\]
若进一步要求
\[
\frac{\pi^2u^2\varphi^{-2(N+1)}}{1-\varphi^{-2}}
\le
\log(1+\varepsilon),
\]
则立得
\[
\frac{C_{\varphi,N}(u)}{C_\varphi(u)}\le 1+\varepsilon.
\]
最后，把上述两条约束分别对 \(N\) 求解，便得到
\[
N\ \ge\
-1+\max\!\left\{
\log_\varphi(2\pi|u|),\
\frac12\log_\varphi\!\left(
\frac{\pi^2u^2}{(1-\varphi^{-2})\log(1+\varepsilon)}
\right)
\right\}.
\]
证毕。
\end{proof}

\noindent
上述七条结果表明，Fibonacci 模零谱与共振整函数 \(\mathcal F_\varphi\) 已经形成一条严格可逆的算术--解析双轨：一方面，谱零点集合 \(\mathcal Z_m\) 不再只是若干候选余类的松散并，而可被零余类覆盖的 flag nerve、约数兼容图的 clique 结构及其加权 Euler 特征精确决定，并进一步把碰撞上界压缩为有限图论计算；另一方面，\(\mathcal F_\varphi\) 的局部解析则不再停留于 Laguerre--P\'olya 性与 Jensen 多项式全实根性这样的抽象实根性，而会在对数 Taylor 系数、倒数零矩、偶系数 Newton 递推与截断深度预算中外显为完全可审计的 \(\zeta(2r)\)--\(\varphi\) 谱矩闭式。于是，零余类几何、有限交 nerve、碰撞上界与共振常数计算深度便被统一收束到同一条由 Fibonacci 约数格与黄金比例谱矩共同支配的闭合链条之中。

\paragraph{奇倍嵌套链、约数格指标容斥、线性零点增量与六倍窗密度指数率}
在零余类族的 \(\gcd\)-刚性交结构与谱零点集合的加权 clique--Euler 公式已经建立之后，还可继续沿约数格方向抽取出下列四条直接后果。它们分别刻画整除链上的奇倍嵌套--偶倍排斥二分律、以相容指标书写的约数格精确容斥公式、pairwise-incompatible 约数族所导致的线性零点增量下界，以及六倍窗上谱零点相对密度的精确指数率。

\begin{theorem}[零余类沿整除链的奇倍嵌套与偶倍排斥]\label{thm:xi-time-part72-zero-coset-divisibility-nested-disjoint-dichotomy}
\leanverified{paper\_xi\_time\_part72\_zero\_coset\_divisibility\_nested\_disjoint\_dichotomy}
设 \(M\) 为偶整数，并令
\[
g\mid h\mid \frac{M}{2}.
\]
则有完全二分律
\[
S_g(M)\subseteq S_h(M)
\iff
\frac{h}{g}\ \text{为奇数},
\]
\[
S_g(M)\cap S_h(M)=\varnothing
\iff
\frac{h}{g}\ \text{为偶数}.
\]
特别地，对任意整除链
\[
g_1\mid g_2\mid \cdots \mid g_r\mid \frac{M}{2},
\]
若所有相邻商
\[
\frac{g_{i+1}}{g_i}\qquad (1\le i<r)
\]
都为奇数，则
\[
S_{g_1}(M)\subseteq S_{g_2}(M)\subseteq \cdots \subseteq S_{g_r}(M),
\qquad
\left|\bigcup_{i=1}^{r}S_{g_i}(M)\right|=g_r.
\]
反之，只要存在某个 \(i\) 使 \(g_{i+1}/g_i\) 为偶数，便有
\[
S_{g_i}(M)\cap S_{g_{i+1}}(M)=\varnothing.
\]
\end{theorem}
\begin{proof}
由定理 \ref{thm:xi-fold-zero-coset-v2-intersection-rigidity}，
\[
S_g(M)\cap S_h(M)=
\begin{cases}
S_{\gcd(g,h)}(M),& v_2(g)=v_2(h),\\[2mm]
\varnothing,& v_2(g)\neq v_2(h).
\end{cases}
\]
在 \(g\mid h\) 的情形，
\[
\gcd(g,h)=g.
\]
又
\[
v_2(g)=v_2(h)
\iff
v_2(h/g)=0
\iff
\frac{h}{g}\ \text{为奇数}.
\]
因此当 \(h/g\) 为奇数时，
\[
S_g(M)\cap S_h(M)=S_g(M),
\]
从而 \(S_g(M)\subseteq S_h(M)\)；当 \(h/g\) 为偶数时，
\[
v_2(g)\neq v_2(h),
\]
故交集为空。

对整除链逐步应用上述二分律，便得到链式包含关系与相邻偶商导致的互斥结论。若全部相邻商皆为奇数，则整条链单调嵌套，故其并集等于最大项 \(S_{g_r}(M)\)，于是
\[
\left|\bigcup_{i=1}^{r}S_{g_i}(M)\right|=|S_{g_r}(M)|=g_r.
\]
证毕。
\end{proof}

\begin{theorem}[谱零点集合的约数格指标容斥公式]\label{thm:xi-time-part72-zero-spectrum-divisor-lattice-indicator-inclusion-exclusion}
\leanverified{paper\_xi\_time\_part72\_zero\_spectrum\_divisor\_lattice\_indicator\_inclusion\_exclusion}
沿用定理 \ref{thm:xi-time-part71-zero-spectrum-weighted-clique-euler} 的记号，令
\[
\mathcal D_m:=
\left\{
d:\ d\mid (m+2),\ d<m+2,\ F_d\mid \frac{F_{m+2}}{2}
\right\}.
\]
对每个非空子集 \(I\subseteq \mathcal D_m\)，定义
\[
\gcd(I):=\gcd\{d:\ d\in I\},
\]
并置相容指标
\[
\mathcal C(I)
\iff
2F_{\gcd(d,e)}\mid (F_e-F_d)
\qquad
(\forall\, d,e\in I,\ d\neq e).
\]
则谱零点集合满足完全显式的约数格容斥展开
\[
|\mathcal Z_m|
=
\sum_{\varnothing\neq I\subseteq \mathcal D_m}
(-1)^{|I|+1}\mathbf 1_{\mathcal C(I)}\,F_{\gcd(I)}.
\]
换言之，\(\mathcal Z_m\) 的基数可由 \((m+2)\) 的约数格与 Fibonacci 差值的整除判别完全决定，而无需再对频率 \(k\) 作逐点扫描。
\end{theorem}
\begin{proof}
由定理 \ref{thm:xi-time-part71-zero-spectrum-weighted-clique-euler}，
\[
|\mathcal Z_m|
=
\sum_{\varnothing\neq C\in \mathrm{Cliq}(\Gamma_m)}
(-1)^{|C|+1}F_{\gcd(C)},
\]
其中 \(\Gamma_m\) 的边关系定义为
\[
d\sim e
\iff
2F_{\gcd(d,e)}\mid (F_e-F_d).
\]
因此，对任意非空子集 \(I\subseteq \mathcal D_m\)，
\[
I\in \mathrm{Cliq}(\Gamma_m)
\iff
\mathcal C(I).
\]
于是把 clique 求和改写为对全部非空子集的求和，并以指标函数 \(\mathbf 1_{\mathcal C(I)}\) 过滤不相容子集，便得到
\[
|\mathcal Z_m|
=
\sum_{\varnothing\neq I\subseteq \mathcal D_m}
(-1)^{|I|+1}\mathbf 1_{\mathcal C(I)}\,F_{\gcd(I)}.
\]
证毕。
\end{proof}

\begin{corollary}[pairwise-incompatible 约数族给出线性零点增量下界]\label{cor:xi-time-part72-zero-spectrum-pairwise-incompatible-linear-lower-bound}
沿用定理 \ref{thm:xi-time-part72-zero-spectrum-divisor-lattice-indicator-inclusion-exclusion} 的记号。若子集
\[
\mathcal A\subseteq \mathcal D_m
\]
满足对任意不同的 \(d,e\in\mathcal A\) 都有
\[
2F_{\gcd(d,e)}\nmid (F_e-F_d),
\]
则相应余类族两两不交，从而
\[
|\mathcal Z_m|
\ge
\sum_{d\in\mathcal A}F_d.
\]
特别地，当 \(m+2\equiv 0\pmod 6\) 且
\[
d:=\frac{m+2}{2},
\]
由推论 \ref{cor:fold-zero-half-index-multiple6} 已知
\[
F_d\mid \frac{F_{m+2}}{2},
\qquad
|\mathcal Z_m|\ge F_d.
\]
若再存在 \(e\in\mathcal D_m\) 满足
\[
2F_{\gcd(d,e)}\nmid (F_e-F_d),
\]
则立刻提升为
\[
|\mathcal Z_m|\ge F_d+F_e.
\]
\end{corollary}
\begin{proof}
由定理 \ref{thm:xi-time-part71-zero-coset-finite-intersection-pairwise-compatibility}，对不同的 \(d,e\in \mathcal A\)，条件
\[
2F_{\gcd(d,e)}\nmid (F_e-F_d)
\]
等价于
\[
S_{F_d}(F_{m+2})\cap S_{F_e}(F_{m+2})=\varnothing.
\]
故族
\[
\{S_{F_d}(F_{m+2}):\ d\in\mathcal A\}
\]
两两不交。于是
\[
\left|
\bigcup_{d\in\mathcal A}S_{F_d}(F_{m+2})
\right|
=
\sum_{d\in\mathcal A}|S_{F_d}(F_{m+2})|
=
\sum_{d\in\mathcal A}F_d,
\]
其中最后一步使用定理 \ref{thm:fold-zero-coset-rigidity} 中
\[
|S_g(M)|=g.
\]
再由推论 \ref{cor:fold-zero-coset-union}，
\[
\bigcup_{d\in\mathcal A}S_{F_d}(F_{m+2})
\subseteq
\mathcal Z_m,
\]
遂得所示下界。

最后一段只需把 \(d=(m+2)/2\) 的六倍窗见证余类与一个与之 pairwise-incompatible 的 \(e\) 一并代入上述结论即可。证毕。
\end{proof}

\begin{theorem}[六倍窗上谱零点相对密度的精确指数率]\label{thm:xi-time-part72-zero-density-exact-exponential-rate}
沿用定理 \ref{thm:xi-time-part70b-dark-band-half-order-scale} 的记号。若
\[
m+2\equiv 0\pmod 6,
\qquad
M:=F_{m+2},
\]
则有极限公式
\[
\lim_{\substack{m\to\infty\\ m+2\equiv 0\!\!\!\pmod 6}}
\frac{1}{m}\log\frac{|\mathcal Z_m|}{F_{m+2}}
=
-\frac12\log\varphi.
\]
等价地，
\[
\frac{|\mathcal Z_m|}{F_{m+2}}
=
\exp\!\left(
-\frac{m}{2}\log\varphi+o(m)
\right)
\]
沿六倍窗子序列成立。
\end{theorem}
\begin{proof}
由定理 \ref{thm:xi-time-part70b-dark-band-half-order-scale}，
\[
\frac{1}{m}\log |\mathcal Z_m|
\longrightarrow
\frac12\log\varphi
\]
沿子序列 \(m+2\equiv 0\pmod 6\) 成立。另一方面，Binet 公式给出
\[
\log F_{m+2}=(m+2)\log\varphi+O(1),
\]
故
\[
\frac{1}{m}\log F_{m+2}\longrightarrow \log\varphi.
\]
两式相减即可得到
\[
\frac{1}{m}\log\frac{|\mathcal Z_m|}{F_{m+2}}
\longrightarrow
\frac12\log\varphi-\log\varphi
=
-\frac12\log\varphi.
\]
等价表述只是把极限式改写为指数渐近。证毕。
\end{proof}

\noindent
上述四条结果表明，零余类几何在有限交截判据与加权 clique--Euler 公式之外，还具有一条更细的约数格层级结构：沿整除链，余类并不会出现任意复杂的部分重叠，而只能在“奇商嵌套”与“偶商互斥”两种情形之间二分；一旦把这一二分律与相容指标容斥结合，任何 pairwise-incompatible 约数簇都会立刻转化为 \(\mathcal Z_m\) 的线性可加增量；而在六倍窗上，这些局部增量又与 half-order dark band 的双边界闭合到同一个精确指数率 \(\varphi^{-m/2}\)。于是，谱零点集合的算术几何便不仅可被精确计数，而且可在约数格上作系统的局部增量构造，并在特定同余窗上获得完全锁定的密度衰减主尺度。

\paragraph{奇商偏序下的零余类压缩生成、唯一极小覆盖与精确计数}
在定理 \ref{thm:xi-time-part71-zero-coset-finite-intersection-pairwise-compatibility} 已给出零余类族任意有限交截的完全分类、定理 \ref{thm:xi-time-part72-zero-coset-divisibility-nested-disjoint-dichotomy} 已把整除链上的交叠压缩为奇倍嵌套与偶倍排斥二分律之后，这条主线还可继续向“表示压缩”推进一步：零点集合 \(\mathcal Z_m\) 并不需要由全部候选余类来生成，而可唯一压缩到奇商偏序下的极大元族；相应的精确计数也可从全约数格容斥进一步压缩为该极大生成族上的包含--排斥闭式。

\begin{theorem}[零余类之间的全局包含关系由奇商整除完全刻画]\label{thm:xi-time-part72a-zero-coset-global-inclusion-odd-divisibility}
\leanpartial{paper\_xi\_time\_part72a\_zero\_coset\_global\_inclusion\_odd\_divisibility}{the requested Lean signature omits the paper's hypothesis that \(M\) is even; at \(M=0\), \(g=1\), and \(h=2\), both zero-cosets equal \(\{0\}\) but \(h/g=2\) is not odd}
设 \(M\) 为偶整数，且 \(g,h\mid M/2\)。则
\[
S_g(M)\subseteq S_h(M)
\iff
g\mid h\ \text{且}\ \frac{h}{g}\ \text{为奇数}.
\]
\end{theorem}
\begin{proof}
若 \(g\mid h\) 且 \(h/g\) 为奇数，则结论已由定理 \ref{thm:xi-time-part72-zero-coset-divisibility-nested-disjoint-dichotomy} 的第一条给出。

反之，设 \(S_g(M)\subseteq S_h(M)\)。则
\[
S_g(M)\cap S_h(M)=S_g(M),
\]
因而交集非空。由定理 \ref{thm:xi-time-part71-zero-coset-finite-intersection-pairwise-compatibility} 在 \(r=2\) 的情形，
\[
\left|S_g(M)\cap S_h(M)\right|=\gcd(g,h).
\]
又由定理 \ref{thm:fold-zero-coset-rigidity}，
\[
|S_g(M)|=g.
\]
于是
\[
\gcd(g,h)=|S_g(M)\cap S_h(M)|=|S_g(M)|=g,
\]
从而 \(g\mid h\)。再由交集非空判据，
\[
2\gcd(g,h)\mid (h-g),
\]
即
\[
2g\mid g\left(\frac{h}{g}-1\right).
\]
故 \(2\mid (h/g-1)\)，也就是 \(h/g\) 为奇数。证毕。
\end{proof}

\begin{theorem}[谱零点集合在奇商偏序下的唯一极小生成族]\label{thm:xi-time-part72a-zero-spectrum-maximal-odd-divisibility-compression}
\leanverified{paper\_xi\_time\_part72a\_zero\_spectrum\_maximal\_odd\_divisibility\_compression}
沿用推论 \ref{cor:fold-zero-coset-union} 与推论 \ref{cor:xi-fold-zero-set-v2-semilattice-decomposition} 的记号。记
\[
M:=F_{m+2},
\qquad
\mathcal D_m:=
\left\{
d:\ d\mid (m+2),\ d<m+2,\ F_d\mid \frac{M}{2}
\right\},
\]
\[
\mathcal G_m:=\{F_d:\ d\in \mathcal D_m\}.
\]
则
\[
\mathcal Z_m=\bigcup_{g\in \mathcal G_m}S_g(M).
\]
在 \(\mathcal G_m\) 上定义偏序
\[
g\preceq h
\iff
S_g(M)\subseteq S_h(M),
\]
等价地，由定理 \ref{thm:xi-time-part72a-zero-coset-global-inclusion-odd-divisibility}，
\[
g\preceq h
\iff
g\mid h\ \text{且}\ \frac{h}{g}\ \text{为奇数}.
\]
记 \(\mathcal M_m\subseteq \mathcal G_m\) 为全部 \(\preceq\)-极大元组成的集合。则成立：
\begin{enumerate}
  \item 谱零点集合具有压缩表示
  \[
  \mathcal Z_m=\bigcup_{g\in \mathcal M_m}S_g(M).
  \]
  \item \(\mathcal M_m\) 是满足上式的唯一极小生成族：若 \(\mathcal A\subseteq \mathcal G_m\) 也满足
  \[
  \mathcal Z_m=\bigcup_{g\in \mathcal A}S_g(M),
  \]
  则必有
  \[
  \mathcal M_m\subseteq \mathcal A.
  \]
\end{enumerate}
\end{theorem}
\begin{proof}
第一条只需删除全部非极大元。事实上，若 \(g\in \mathcal G_m\setminus \mathcal M_m\)，则存在 \(h\in \mathcal G_m\) 使 \(g\prec h\)。由定理 \ref{thm:xi-time-part72a-zero-coset-global-inclusion-odd-divisibility}，
\[
S_g(M)\subseteq S_h(M),
\]
故从并集中去掉 \(S_g(M)\) 不改变并集。对所有非极大元依次执行此操作，即得
\[
\mathcal Z_m=\bigcup_{g\in \mathcal M_m}S_g(M).
\]

现证唯一极小性。任取 \(g\in \mathcal M_m\)。若 \(h\in \mathcal G_m\setminus\{g\}\) 且
\[
S_g(M)\cap S_h(M)\neq \varnothing,
\]
则由定理 \ref{thm:xi-time-part71-zero-coset-finite-intersection-pairwise-compatibility}，
\[
S_g(M)\cap S_h(M)=S_{\gcd(g,h)}(M)
\]
且
\[
\gcd(g,h)<g,
\]
否则若 \(\gcd(g,h)=g\)，则 \(g\mid h\)，再由交非空判据可得 \(h/g\) 为奇数，从而 \(g\prec h\)，与 \(g\) 的极大性矛盾。

取 \(S_g(M)\) 的标准代表点
\[
k_g:=\frac{M}{2g}\in S_g(M).
\]
若 \(k_g\in S_h(M)\) 对某个 \(h\neq g\) 成立，则由上式知
\[
k_g\in S_g(M)\cap S_h(M)=S_{\gcd(g,h)}(M)
\]
且 \(d:=\gcd(g,h)<g\)。由于交集非空，\(g/d\) 必为奇数。于是 \(S_d(M)\subseteq S_g(M)\)，并且按定理 \ref{thm:fold-zero-coset-rigidity} 的显式描述，\(S_d(M)\) 在 \(S_g(M)\) 内对应于模 \(g/d\) 的单一余类
\[
t\equiv \frac{g/d-1}{2}\pmod{g/d}.
\]
而点 \(k_g\) 对应于参数 \(t=0\)。因为 \(g/d>1\) 且为奇数，
\[
\frac{g/d-1}{2}\not\equiv 0\pmod{g/d},
\]
故 \(k_g\notin S_d(M)\)，矛盾。

因此
\[
k_g\notin \bigcup_{h\in \mathcal G_m\setminus\{g\}}S_h(M),
\]
从而
\[
S_g(M)\nsubseteq \bigcup_{h\in \mathcal G_m\setminus\{g\}}S_h(M).
\]
于是任何生成整个 \(\mathcal Z_m\) 的子族都必须包含 \(g\)。由于 \(g\in \mathcal M_m\) 任意，遂得
\[
\mathcal M_m\subseteq \mathcal A.
\]
这就证明了唯一极小性。证毕。
\end{proof}

\begin{theorem}[极大奇商生成族上的谱零点集合精确容斥公式]\label{thm:xi-time-part72a-zero-spectrum-maximal-family-inclusion-exclusion}
沿用定理 \ref{thm:xi-time-part72a-zero-spectrum-maximal-odd-divisibility-compression} 的记号，写
\[
\mathcal M_m=\{g_1,\dots,g_r\}.
\]
则谱零点集合的基数满足精确公式
\[
|\mathcal Z_m|
=
\sum_{\varnothing\neq I\subseteq \{1,\dots,r\}}
(-1)^{|I|+1}
\mathbf 1_{\mathrm{comp}(I)}
\gcd(g_i:\ i\in I),
\]
其中相容指标 \(\mathbf 1_{\mathrm{comp}(I)}\) 定义为：当且仅当对全部 \(i,j\in I\) 都有
\[
2\gcd(g_i,g_j)\mid (g_j-g_i)
\]
时，\(\mathbf 1_{\mathrm{comp}(I)}=1\)；否则 \(\mathbf 1_{\mathrm{comp}(I)}=0\)。
\end{theorem}
\begin{proof}
由定理 \ref{thm:xi-time-part72a-zero-spectrum-maximal-odd-divisibility-compression}，
\[
\mathcal Z_m=\bigcup_{\nu=1}^{r}S_{g_\nu}(M).
\]
对该有限并应用包含--排斥，
\[
|\mathcal Z_m|
=
\sum_{\varnothing\neq I\subseteq \{1,\dots,r\}}
(-1)^{|I|+1}
\left|
\bigcap_{i\in I}S_{g_i}(M)
\right|.
\]
再由定理 \ref{thm:xi-time-part71-zero-coset-finite-intersection-pairwise-compatibility}，交截非空当且仅当上述两两相容条件成立，而一旦非空就有
\[
\left|
\bigcap_{i\in I}S_{g_i}(M)
\right|
=
\gcd(g_i:\ i\in I).
\]
代回即得所示公式。证毕。
\end{proof}

\begin{corollary}[极大生成族 pairwise-incompatible 时的精确直和公式]\label{cor:xi-time-part72a-zero-spectrum-maximal-family-pairwise-incompatible-exact-sum}
沿用定理 \ref{thm:xi-time-part72a-zero-spectrum-maximal-odd-divisibility-compression} 的记号。若极大生成族 \(\mathcal M_m=\{g_1,\dots,g_r\}\) 满足任意两元都不相容，即
\[
2\gcd(g_i,g_j)\nmid (g_j-g_i)
\qquad (i\neq j),
\]
则对应余类两两不交，从而
\[
|\mathcal Z_m|=\sum_{\nu=1}^{r}g_\nu.
\]
\end{corollary}
\begin{proof}
由定理 \ref{thm:xi-time-part71-zero-coset-finite-intersection-pairwise-compatibility}，上述条件等价于
\[
S_{g_i}(M)\cap S_{g_j}(M)=\varnothing
\qquad (i\neq j).
\]
故
\[
\mathcal Z_m=\bigsqcup_{\nu=1}^{r}S_{g_\nu}(M).
\]
再由定理 \ref{thm:fold-zero-coset-rigidity} 的 \(|S_g(M)|=g\)，可得
\[
|\mathcal Z_m|
=
\sum_{\nu=1}^{r}|S_{g_\nu}(M)|
=
\sum_{\nu=1}^{r}g_\nu.
\]
证毕。
\end{proof}

\noindent
上述结果表明，零余类族的算术几何还具有一层此前未被完全显化的“压缩表示”刚性：有限交截虽由全体候选余类的 pairwise 相容性控制，但真正生成 \(\mathcal Z_m\) 的并不需要整张约数图，而只需要奇商偏序下的极大元族；这一压缩不仅是存在性的，而且是唯一极小的。相应地，谱零点计数也可从全约数格的指标容斥，进一步压缩为极大生成族上的精确包含--排斥；而在 pairwise-incompatible 的情形下，全部交叠复杂性彻底消失，精确计数退化为一条纯粹的 Fibonacci 权重直和公式。于是，零余类交叠、包含、压缩表示与 exact counting 四者就在同一条 odd-divisibility 主线上获得了闭合。

\paragraph{pairwise-incompatible 零块族的同步信息削减}
在推论 \ref{cor:xi-time-part72-zero-spectrum-pairwise-incompatible-linear-lower-bound} 已把 pairwise-incompatible 约数族转化为 \(|\mathcal Z_m|\) 的线性可加增量，而定理 \ref{thm:xi-fold-zero-dyadic-tower-synchronous-information-gap-improvement} 已在 dyadic 塔特例下把这一增量同步传递到碰撞、KL、峰值与最大纤维偏离之后，可以看到：这一同步机制本身并不依赖 dyadic 结构，而只依赖 pairwise incompatibility 所保证的零块互不相交。

\begin{theorem}[pairwise-incompatible Fibonacci 约数族触发的同步信息削减]\label{thm:xi-time-part72b-pairwise-incompatible-synchronous-information-gap}
\leanverified{paper\_xi\_time\_part72b\_pairwise\_incompatible\_synchronous\_information\_gap}
沿用定理 \ref{thm:xi-time-part72-zero-spectrum-divisor-lattice-indicator-inclusion-exclusion} 的记号，并沿用定理 \ref{thm:fold-zero-uncertainty} 的记号。设
\[
\mathcal A\subseteq \mathcal D_m
\]
满足对任意不同的 \(d,e\in\mathcal A\)，都有
\[
2F_{\gcd(d,e)}\nmid (F_e-F_d).
\]
记
\[
L_{\mathcal A}:=\sum_{d\in\mathcal A}F_d.
\]
则
\[
|\mathcal Z_m|\ge L_{\mathcal A}.
\]
因而同时成立
\[
\mathsf{Col}_m
\le
1-\frac{L_{\mathcal A}}{F_{m+2}},
\qquad
D_{\mathrm{KL}}(p_m\|u_m)
\le
\log\!\bigl(F_{m+2}-L_{\mathcal A}\bigr),
\]
\[
S_m\le 2^m\sqrt{F_{m+2}-1-L_{\mathcal A}},
\]
\[
M_m-\bar d_m
\le
2^m\sqrt{\frac{F_{m+2}-1-L_{\mathcal A}}{F_{m+2}}}.
\]
\end{theorem}
\begin{proof}
由推论 \ref{cor:xi-time-part72-zero-spectrum-pairwise-incompatible-linear-lower-bound}，
\[
|\mathcal Z_m|\ge \sum_{d\in\mathcal A}F_d=L_{\mathcal A}.
\]
将该下界分别代入定理 \ref{thm:fold-collision-zero-reduction} 与定理
\ref{thm:fold-zero-uncertainty} 的第 \(2\)、\(3\)、\(4\) 条，便得到所示碰撞上界、KL 上界、Fourier 峰值上界以及最大纤维偏离上界。证毕。
\end{proof}

\noindent
因此，定理 \ref{thm:xi-fold-zero-dyadic-tower-synchronous-information-gap-improvement} 只是这一一般机制在 dyadic 塔族上的特例。真正驱动同步削减的并非 \(v_2\)-塔本身，而是 Fibonacci 零块之间的 pairwise incompatibility：只要若干约数对应的零块能够同时并入 \(\mathcal Z_m\) 而不发生交叠，它们所贡献的零点增量便会无损地传递到碰撞、相对熵、峰值振幅与最大纤维偏离这四个量上。

\endinput


\paragraph{二维原子层、非刚性迹系统、激活零尺度根与零温四态核心的三尺度分裂}
在真实输入 \(40\) 状态核的二维加权 \(\zeta\) 分母闭式、一参数输出位势的激活分支，以及零温四态地基层已经建立之后，还可把原子层、非刚性残余谱、局部激活层与零温核心层进一步压缩为下列若干条彼此衔接的结构后果。前两条仍对应对原子因子的中心投影式相位盲读出；紧随其后的两条把去除 rigid branches 后的全部非刚性迹压缩为单一八次 core 因子的 Newton--Cayley 闭系统，并进一步抽出 \(v=1\) 超曲面与 \((u,v)=(0,1)\) 尖点处的谱秩塌缩；随后的几条则转入相位敏感的纤维内读出，分别刻画 \(u=1\) 处的极点阶、Abel 有限部平移、激活分支的横截性、去支后的局部无碰撞窗口以及正实轴上的局部谱壁垒；最后再把局部奇性独占与零温 Perron 幅度统一到同一条谱分裂链中。

\begin{theorem}[二维加权 \texorpdfstring{$\zeta$}{zeta} 的原子因子在 ghost 层对碰撞权完全失明]\label{thm:xi-time-part73-twoparam-atomic-ghost-u-blindness}
\leanverified{paper\_xi\_time\_part73\_twoparam\_atomic\_ghost\_u\_blindness}
沿用定理 \ref{thm:real-input-40-det-uv-2plus8} 与推论 \ref{cor:real-input-40-zeta-uv-sqrtv-eigs} 的记号，并记
\[
\zeta(z;u,v)
=
\frac{1}{\Delta(z;u,v)}
=
\frac{1}{1-vz^2}\,\zeta_{\mathrm{core}}(z;u,v),
\qquad
\zeta_{\mathrm{core}}(z;u,v):=-P_8(z;u,v)^{-1}.
\]
定义原子 ghost 系数 \(a_n^{\mathrm{atom}}(u,v)\) 为
\[
\log\frac{1}{1-vz^2}
=
\sum_{n\ge 1}\frac{a_n^{\mathrm{atom}}(u,v)}{n}z^n.
\]
则对任意 \(n\ge 1\) 都有精确公式
\[
\boxed{\;
a_n^{\mathrm{atom}}(u,v)
=
\begin{cases}
2v^{n/2},& 2\mid n,\\[4pt]
0,& 2\nmid n,
\end{cases}\;}
\]
因而该原子层对碰撞参数 \(u\) 完全独立。若进一步取 \(v>0\)，则同一公式可重写为
\[
\boxed{\;
a_n^{\mathrm{atom}}(u,v)
=
(\sqrt v)^n+(-\sqrt v)^n.\;}
\]
特别地，所有奇阶 ghost/trace 审计都对因子 \((1-vz^2)^{-1}\) 完全失明：对任意 \(m\ge 0\)，其对
\[
\operatorname{Tr}(M_{u,v}^{\,2m+1})
\]
的原子贡献恒为 \(0\)。
\end{theorem}
\begin{proof}
由
\[
-\log(1-vz^2)=\sum_{k\ge 1}\frac{v^k}{k}z^{2k}
\]
并与 ghost 记号
\[
\log\frac{1}{1-vz^2}
=
\sum_{n\ge 1}\frac{a_n^{\mathrm{atom}}(u,v)}{n}z^n
\]
逐项比较系数，立得
\[
a_{2k}^{\mathrm{atom}}(u,v)=2v^k,
\qquad
a_{2k+1}^{\mathrm{atom}}(u,v)=0.
\]
这已经给出对 \(u\) 的完全独立性。

当 \(v>0\) 时，推论 \ref{cor:real-input-40-zeta-uv-sqrtv-eigs} 表明因子 \(1-vz^2\) 在谱上对应两条 rigid branches
\[
\lambda_\pm=\pm\sqrt v,
\]
故其对 \(n\)-步迹的贡献正是
\[
(\sqrt v)^n+(-\sqrt v)^n,
\]
而这与上式完全一致。于是奇数阶贡献必为零，从而任何只读取奇阶 ghost/trace 的审计函数都无法检测该原子因子。证毕。
\end{proof}

\begin{theorem}[去除 rigid branches 后的非刚性迹由八次 core 因子完全编码]\label{thm:xi-time-part73-nonrigid-trace-newton-cayley-closure}
沿用定理 \ref{thm:real-input-40-det-uv-2plus8} 与推论 \ref{cor:real-input-40-zeta-uv-sqrtv-eigs} 的记号。定义正规化 core 多项式
\[
\mathcal Q_8(z;u,v):=-P_8(z;u,v),
\qquad
\mathcal Q_8(0;u,v)=1.
\]
并写成
\[
\mathcal Q_8(z;u,v)=1+c_1z+\cdots+c_8z^8,
\]
其中若 \(\deg_z\mathcal Q_8<8\)，则约定高阶系数 \(c_j\) 自动取 \(0\)。当 \(v>0\) 时，再定义非刚性迹
\[
R_n(u,v):=\operatorname{Tr}(M_{u,v}^n)-(\sqrt v)^n-(-\sqrt v)^n
\qquad(n\ge 1).
\]
则成立：
\begin{enumerate}
  \item 生成函数恒等式
  \[
  \boxed{\;
  -\frac{z\,\partial_z\mathcal Q_8(z;u,v)}{\mathcal Q_8(z;u,v)}
  =
  \sum_{n\ge 1}R_n(u,v)\,z^n.\;}
  \]
  \item 当 \(1\le n\le 8\) 时，
  \[
  \boxed{\;
  R_n+c_1R_{n-1}+\cdots+c_{n-1}R_1+n\,c_n=0.\;}
  \]
  \item 当 \(n\ge 9\) 时，
  \[
  \boxed{\;
  R_n+c_1R_{n-1}+c_2R_{n-2}+\cdots+c_8R_{n-8}=0.\;}
  \]
\end{enumerate}
因此，二维加权核除去刚性对 \(\pm\sqrt v\) 之后的全部非刚性迹谱，完全由单一八次多项式 \(\mathcal Q_8\) 编码，并形成阶数至多为 \(8\) 的闭线性递推系统。
\end{theorem}
\begin{proof}
由定理 \ref{thm:real-input-40-det-uv-2plus8}，
\[
\Delta(z;u,v)=(1-vz^2)\mathcal Q_8(z;u,v).
\]
再由推论 \ref{cor:real-input-40-zeta-uv-sqrtv-eigs}，当 \(v>0\) 时矩阵 \(M_{u,v}\) 的两条 rigid branches 恰为
\[
\lambda_\pm=\pm\sqrt v.
\]
因此存在至多 \(8\) 个非刚性非零特征值
\[
\eta_1,\dots,\eta_d
\qquad(d\le 8),
\]
使得
\[
\mathcal Q_8(z;u,v)=\prod_{j=1}^{d}(1-z\eta_j),
\qquad
R_n(u,v)=\sum_{j=1}^{d}\eta_j^n.
\]
于是
\[
-\frac{z\,\partial_z\mathcal Q_8(z;u,v)}{\mathcal Q_8(z;u,v)}
=
\sum_{j=1}^{d}\frac{z\eta_j}{1-z\eta_j}
=
\sum_{j=1}^{d}\sum_{n\ge 1}\eta_j^n z^n
=
\sum_{n\ge 1}R_n(u,v)\,z^n,
\]
得到第一条。

其后两条只是标准 Newton--Girard 恒等式在
\[
\mathcal Q_8(z;u,v)=1+c_1z+\cdots+c_8z^8
\]
与幂和 \(R_n=\sum_j\eta_j^n\) 上的专门化：对 \(1\le n\le 8\) 得到非齐次关系
\[
R_n+c_1R_{n-1}+\cdots+c_{n-1}R_1+n\,c_n=0,
\]
而对 \(n\ge 9\) 则转为齐次关系
\[
R_n+c_1R_{n-1}+\cdots+c_8R_{n-8}=0.
\]
若 \(d<8\)，则高阶系数 \(c_{d+1},\dots,c_8\) 本来即为零，故同一公式仍然成立。证毕。
\end{proof}

\begin{theorem}[二维加权核在 \texorpdfstring{$v=1$}{v=1} 上发生共维一谱秩塌缩]\label{thm:xi-time-part73-v1-rank-collapse-codim2-tip}
沿用定理 \ref{thm:xi-time-part73-nonrigid-trace-newton-cayley-closure} 的记号。则有：
\begin{enumerate}
  \item 对任意 \(u\)，在超曲面 \(v=1\) 上都成立
  \[
  \deg_z \Delta(z;u,1)\le 9,
  \qquad
  \dim\ker M_{u,1}\ge 31.
  \]
  \item 在共维二点 \((u,v)=(0,1)\) 上，显式有
  \[
  P_8(z;0,1)=-2z^5-6z^4-z^3+5z^2+z-1,
  \]
  从而
  \[
  \deg_z \Delta(z;0,1)\le 7,
  \qquad
  \dim\ker M_{0,1}\ge 33.
  \]
  \item 相应地，上一条定理中的非刚性迹递推在 \(v=1\) 上自动降为阶数至多 \(7\) 的系统，而在 \((u,v)=(0,1)\) 处进一步降为阶数至多 \(5\) 的系统。
\end{enumerate}
\end{theorem}
\begin{proof}
由定理 \ref{thm:real-input-40-det-uv-2plus8}，
\[
P_8(z;u,v)
=(u^2v^4-u^2v^3)z^8+\cdots
=u^2v^3(v-1)z^8+\cdots.
\]
因此当 \(v=1\) 时，\(z^8\) 项恒消失，故
\[
\deg_z P_8(z;u,1)\le 7.
\]
再由
\[
\Delta(z;u,1)=(z^2-1)P_8(z;u,1),
\]
便得到
\[
\deg_z\Delta(z;u,1)\le 2+7=9.
\]
另一方面，\(\Delta(z;u,v)=\det(I-zM_{u,v})\) 的 \(z\)-次数恰等于 \(M_{u,v}\) 的全部非零特征值总数（按代数重数计）。由于 \(M_{u,1}\) 为 \(40\times 40\) 矩阵，这意味着
\[
\dim\ker M_{u,1}\ge 40-9=31.
\]

若进一步取 \((u,v)=(0,1)\)，则把 \(u=0,v=1\) 直接代入定理 \ref{thm:real-input-40-det-uv-2plus8} 的显式八次式可得
\[
P_8(z;0,1)=-2z^5-6z^4-z^3+5z^2+z-1,
\]
故
\[
\deg_z P_8(z;0,1)\le 5,
\qquad
\deg_z \Delta(z;0,1)\le 2+5=7.
\]
于是
\[
\dim\ker M_{0,1}\ge 40-7=33.
\]

最后，由定理 \ref{thm:xi-time-part73-nonrigid-trace-newton-cayley-closure}，非刚性迹递推的阶数恰由多项式 \(\mathcal Q_8=-P_8\) 的次数控制。故在 \(v=1\) 上其阶数至多为 \(7\)，而在 \((u,v)=(0,1)\) 处至多为 \(5\)。证毕。
\end{proof}

\endinput


\begin{theorem}[二维加权 \texorpdfstring{$\zeta$}{zeta} 在碰撞零温缩放下坍缩为单一原子模型]\label{thm:xi-time-part73-twoparam-collision-freezing-single-atom-collapse}
沿用定理 \ref{thm:real-input-40-det-uv-2plus8} 的记号。固定 \(w>0\) 与 \(v>0\)，并令
\[
z=\sqrt{\frac{w}{u}},
\qquad
u\to+\infty.
\]
则有极限
\[
\boxed{\;
\lim_{u\to+\infty}\Delta\!\left(\sqrt{\frac{w}{u}};u,v\right)
=
1-vw.\;}
\]
更精确地，
\[
\boxed{\;
P_8\!\left(\sqrt{\frac{w}{u}};u,v\right)
=
-1+vw+O(u^{-1/2}),\;}
\]
从而
\[
\boxed{\;
\Delta\!\left(\sqrt{\frac{w}{u}};u,v\right)
=
1-vw+O(u^{-1/2}).\;}
\]
因此只要 \(vw\neq 1\)，便等价地有
\[
\boxed{\;
\zeta\!\left(\sqrt{\frac{w}{u}};u,v\right)
\longrightarrow
\frac{1}{1-vw}.\;}
\]
换言之，在碰撞零温缩放下，二维 \((u,v)\)-联合势的核心部分完全退化为长度 \(2\)、权重 \(v\) 的单原子 Euler 模型。
\end{theorem}
\begin{proof}
将定理 \ref{thm:real-input-40-det-uv-2plus8} 中 \(P_8(z;u,v)\) 的显式八次式代入
\[
z=\sqrt{\frac{w}{u}}.
\]
逐项估计可见：常数项给出 \(-1\)，而 \(z^2\) 项中的主贡献来自
\[
(uv+5v)z^2=vw+5vw\,u^{-1}.
\]
其余各项都至少带有 \(u^{-1/2}\) 的衰减；更具体地，\(z,z^3,z^5,z^7\) 项为 \(O(u^{-1/2})\)，而 \(z^4,z^6,z^8\) 项为 \(O(u^{-1})\)。故
\[
P_8\!\left(\sqrt{\frac{w}{u}};u,v\right)
=
-1+vw+O(u^{-1/2}).
\]

另一方面，
\[
vz^2-1
=
\frac{vw}{u}-1
=
-1+O(u^{-1}).
\]
于是
\[
\Delta\!\left(\sqrt{\frac{w}{u}};u,v\right)
=
\Bigl(-1+O(u^{-1})\Bigr)
\Bigl(-1+vw+O(u^{-1/2})\Bigr)
=
1-vw+O(u^{-1/2}),
\]
从而得到第一条极限。若 \(vw\neq 1\)，则极限不为零，故取倒数即得 \(\zeta\) 的极限公式。证毕。
\end{proof}

\begin{theorem}[\texorpdfstring{$u=1$}{u=1} 处 \texorpdfstring{$\zeta$}{zeta} 的双--三极点律]\label{thm:xi-time-part73-u1-zeta-double-triple-pole}
\leanverified{paper\_xi\_time\_part73\_u1\_zeta\_double\_triple\_pole}
沿用命题 \ref{prop:real-input-40-output-spectral-collisions} 的记号，并记
\[
\zeta(z,1):=\Delta(z,1)^{-1},
\qquad
\zeta_{\mathrm{core}}(z,1):=Q(z,1)^{-1}.
\]
则有精确分解
\[
\Delta(z,1)=-(z-1)^2(z+1)^3(z^2-3z+1)(z^2-z-1).
\]
因此 \(\zeta(z,1)\) 在 \(z=1\) 具有二阶极点，在 \(z=-1\) 具有三阶极点。去除原子 Witt 因子 \((1-z^2)^{-1}\) 之后，核心 \(\zeta_{\mathrm{core}}(z,1)\) 在 \(z=1\)、\(z=-1\) 分别仅余一阶与二阶极点。
\end{theorem}
\begin{proof}
由命题 \ref{prop:real-input-40-output-spectral-collisions}，
\[
Q(z,1)=(z-1)(z+1)^2(z^2-3z+1)(z^2-z-1).
\]
另一方面，
\[
\Delta(z,1)=(1-z^2)Q(z,1)=-(z-1)(z+1)Q(z,1),
\]
将两式相乘即得所示因式分解。于是 \(\zeta(z,1)=\Delta(z,1)^{-1}\) 在 \(z=1\)、\(z=-1\) 的极点阶分别为 \(2\) 与 \(3\)。而
\[
\zeta_{\mathrm{core}}(z,1)=Q(z,1)^{-1}
\]
只除去原子因子 \((1-z^2)^{-1}\)，故两处极点阶分别降低 \(1\)。证毕。
\end{proof}

\begin{theorem}[原子二周期对 Abel 有限部的精确平移]\label{thm:xi-time-part73-two-cycle-abel-finitepart-shift}
沿用定理 \ref{thm:xi-time-part65-atomic-primitive-abel-shift} 的记号。在 \(u=1\) 的同一核上，设 \(z_\star=\varphi^{-2}\)。则 primitive--core 分解满足
\[
P_{\mathrm{core}}(z)=P(z)-z^2,
\qquad
\log\mathfrak M=\log\mathfrak M_{\mathrm{core}}+z_\star^2.
\]
因此
\[
\log\mathfrak M-\log\mathfrak M_{\mathrm{core}}
=
\varphi^{-4}
=
\frac{7-3\sqrt5}{2},
\]
亦即
\[
\mathfrak M_{\mathrm{core}}=e^{-\varphi^{-4}}\mathfrak M.
\]
\end{theorem}
\begin{proof}
由定理 \ref{thm:xi-time-part65-atomic-primitive-abel-shift}，
\[
P(z,u)=uz^2+P_{\mathrm{core}}(z,u).
\]
在 \(u=1\) 处专门化即得
\[
P_{\mathrm{core}}(z)=P(z)-z^2.
\]
同一定理并给出
\[
\log\mathfrak M=\log\mathfrak M_{\mathrm{core}}+z_\star^2.
\]
又在本核上 \(z_\star=\varphi^{-2}\)，故
\[
z_\star^2=\varphi^{-4}=\frac{7-3\sqrt5}{2}.
\]
代回并指数化，便得到所示 Abel 有限部平移公式。证毕。
\end{proof}

\begin{theorem}[激活分支的横截性与二次项消失]\label{thm:xi-time-part73-activated-branch-unique-zero-scale-singularity}
沿用命题 \ref{prop:real-input-40-activated-branch-linear}、注 \ref{rem:real-input-40-activated-branch-series} 与定理 \ref{thm:xi-output-potential-activated-branch-rigidity} 的记号，并记
\[
G(\lambda,u):=\lambda^8F(\lambda^{-1},u).
\]
则在 \(u=1\) 的某邻域内存在唯一解析分支 \(\lambda_{\mathrm{act}}(u)\) 满足
\[
G(\lambda_{\mathrm{act}}(u),u)=0,
\qquad
\lambda_{\mathrm{act}}(1)=0,
\]
并且其 Taylor 展开满足
\[
\boxed{\;
\lambda_{\mathrm{act}}(u)
=(u-1)+3(u-1)^3-(u-1)^4+18(u-1)^5-22(u-1)^6
+O((u-1)^7).\;}
\]
特别地，二次项严格消失，且
\[
\lambda'_{\mathrm{act}}(1)=1,
\]
从而该分支穿过 \(\lambda=0\) 的方式是横截而非抛物切触。进一步，存在常数 \(\rho>0\) 与 \(\varepsilon>0\)，使得对一切满足 \(|u-1|<\varepsilon\) 的实参数，都有
\[
\boxed{\;
|\lambda_j(u)|\ge \rho
\qquad
(2\le j\le 8).\;}
\]
\end{theorem}
\begin{proof}
由命题 \ref{prop:real-input-40-activated-branch-linear}，
\[
\partial_\lambda G(0,1)=1\neq 0,
\qquad
\lambda_{\mathrm{act}}(1)=0.
\]
故由隐函数定理，\(G(\lambda,u)=0\) 在 \((0,1)\) 邻域内确定唯一解析分支 \(\lambda_{\mathrm{act}}(u)\)，且
\[
\lambda'_{\mathrm{act}}(1)=1.
\]
另一方面，注 \ref{rem:real-input-40-activated-branch-series} 给出
\[
\lambda_{\mathrm{act}}(u)
=(u-1)+3(u-1)^3-(u-1)^4+18(u-1)^5-22(u-1)^6
+O((u-1)^7),
\]
故二次项系数严格为 \(0\)，而线性项系数为 \(1\)，从而该分支以横截方式穿过 \(\lambda=0\)。

最后，定理 \ref{thm:xi-output-potential-activated-branch-rigidity} 已证明：在 \((\lambda,u)=(0,1)\) 的某邻域内，唯一的小根就是 \(\lambda_{\mathrm{act}}(u)\)，其余七条根分支全部远离 \(0\)。因此存在 \(r>0\) 与 \(\varepsilon_0>0\)，使得当 \(|u-1|<\varepsilon_0\) 时，除 \(\lambda_{\mathrm{act}}(u)\) 外的所有根都落在 \(|\lambda|\ge r\) 上。取 \(\rho:=r\)、\(\varepsilon:=\varepsilon_0\) 即得最后一条。证毕。
\end{proof}

\begin{theorem}[激活分支的三次归一化刚性与去支实谱的无碰撞窗口]\label{thm:xi-time-part73-activated-branch-deflation-window}
沿用定理 \ref{thm:xi-time-part73-activated-branch-unique-zero-scale-singularity} 的记号，并记
\[
I_\ast:=
\bigl(0.93283726248596492776\ldots,\ 1.0547964213679192275\ldots\bigr).
\]
定义三次归一化 defect germ
\[
\Xi(u):=\frac{\lambda_{\mathrm{act}}(u)-(u-1)}{(u-1)^3}.
\]
又定义去支因子
\[
\widetilde G(\lambda,u):=\frac{G(\lambda,u)}{\lambda-\lambda_{\mathrm{act}}(u)}.
\]
则成立：
\begin{enumerate}
  \item \(\Xi(u)\) 在 \(u=1\) 处可解析延拓，并满足
  \[
  \Xi(1)=3.
  \]
  \item \(\widetilde G\) 在 \((\lambda,u)=(0,1)\) 的某邻域内解析；
  \item 对每个实参数 \(u\in I_\ast\setminus\{1\}\)，多项式 \(\widetilde G(\cdot,u)\) 的全部实根都是单根。换言之，去除激活分支之后，审计窗口 \(I_\ast\) 内的剩余实谱不再发生双根碰撞。
\end{enumerate}
\end{theorem}
\begin{proof}
由注 \ref{rem:real-input-40-activated-branch-series}，
\[
\lambda_{\mathrm{act}}(u)-(u-1)
=
3(u-1)^3-(u-1)^4+18(u-1)^5-22(u-1)^6+O((u-1)^7).
\]
故
\[
\Xi(u)
=
3-(u-1)+18(u-1)^2-22(u-1)^3+O((u-1)^4),
\]
从而 \(\Xi\) 在 \(u=1\) 处具有可去奇性，并且
\[
\Xi(1)=3.
\]

第二条正是定理 \ref{thm:xi-output-potential-activated-branch-rigidity} 的 Weierstrass 因除结论。

最后，设某个 \(u_0\in I_\ast\setminus\{1\}\) 使 \(\widetilde G(\cdot,u_0)\) 存在实双根 \(\lambda_0\)。由于
\[
G(0,u)=u^3(1-u),
\]
且 \(u_0\neq 1\)，故 \(\lambda_0\neq 0\)。记
\[
z_0:=\lambda_0^{-1}\in \RR.
\]
在 \(\lambda\neq 0\) 上，变换 \(z=\lambda^{-1}\) 为解析双射，而
\[
G(\lambda,u)=\lambda^8F(\lambda^{-1},u)
\]
只差一个处处非零因子 \(\lambda^8\)，故 \(\lambda_0\) 为 \(\widetilde G(\cdot,u_0)\) 的实双根会强迫 \(z_0\) 成为 core 因子 \(F(z,u_0)\) 的实双根；沿一参数输出位势的记号，这等价于 \(Q(z,u_0)\) 在实轴上出现重根。可这与注 \ref{rem:real-input-40-output-potential-degeneracy-phase} 矛盾：该注指出在正实轴上除 \(u=1\) 外，全部退化点恰为
\[
0.008460\ldots,\ 0.745534\ldots,\ 0.932837\ldots,\ 1.054796\ldots,\ 2.090392\ldots,\ 15.522820\ldots,
\]
因而在开区间 \(I_\ast\setminus\{1\}\) 内不存在任何重根参数。故 \(\widetilde G(\cdot,u)\) 的实根在该窗口内全为单根。证毕。
\end{proof}

\begin{corollary}[围绕 \texorpdfstring{$u=1$}{u=1} 的局部谱病态完全由唯一激活小根承担]\label{cor:xi-time-part73-single-activated-branch-all-local-illconditioning}
沿用定理 \ref{thm:xi-time-part73-activated-branch-unique-zero-scale-singularity} 与定理 \ref{thm:xi-time-part73-activated-branch-deflation-window} 的记号。则存在 \(u=1\) 的邻域 \(\mathcal U\) 与在 \((\lambda,u)=(0,1)\) 某邻域内解析的函数 \(\widetilde G(\lambda,u)\)，使得
\[
G(\lambda,u)=\bigl(\lambda-\lambda_{\mathrm{act}}(u)\bigr)\widetilde G(\lambda,u)
\qquad(u\in\mathcal U).
\]
并且存在常数 \(c>0\)，使得对每个 \(u\in\mathcal U\)，\(\widetilde G(\cdot,u)\) 的任一零点 \(\mu(u)\) 都满足
\[
|\mu(u)|\ge c.
\]
因此，\(u=1\) 邻域内全部由小根诱导的局部谱病态都由唯一分支 \(\lambda_{\mathrm{act}}(u)\) 承担；去除该分支之后，其余七支统一恢复到 \(O(1)\) 尺度。
\end{corollary}
\begin{proof}
定理 \ref{thm:xi-time-part73-activated-branch-deflation-window} 的第二条给出
\[
\widetilde G(\lambda,u):=\frac{G(\lambda,u)}{\lambda-\lambda_{\mathrm{act}}(u)}
\]
在 \((\lambda,u)=(0,1)\) 的某邻域内解析，故所示 Weierstrass 因除成立。

另一方面，定理 \ref{thm:xi-time-part73-activated-branch-unique-zero-scale-singularity} 已证明：在 \(u=1\) 的某邻域内，除 \(\lambda_{\mathrm{act}}(u)\) 外，其余全部根都与 \(0\) 保持一致正距离。于是可取某个常数 \(c>0\)，使 \(\widetilde G(\cdot,u)\) 的所有零点统一满足
\[
|\mu(u)|\ge c.
\]
这恰表明全部局部小尺度病态只由 \(\lambda_{\mathrm{act}}(u)\) 一支承担，而去支后的剩余谱支都停留在 \(O(1)\) 距离之外。证毕。
\end{proof}

\begin{theorem}[正实轴上的谱“护城河”]\label{thm:xi-time-part73-positive-real-spectral-moat}
仍在真实输入 \(40\) 状态核的输出位势加权下，取
\[
\Delta(z,u)=(1-uz^2)Q(z,u),
\qquad
u>0.
\]
则正实参数区间
\[
1<u<4
\]
构成一个严格的谱护城河：在该区间内同时不存在显式原子分支 \(z=\pm u^{-1/2}\) 与 \(Q\)-谱的碰撞，也不存在 \(Q\)-谱在整数点 \(z=\pm1\) 的共振。更精确地，
\[
Q(\pm u^{-1/2},u)=0
\iff
u=1,
\]
且
\[
Q(1,u)=-u(u-4)(u-1)(u+1),
\qquad
Q(-1,u)=-(u-1)^2(u^2+6u-2).
\]
故在 \(u>1\) 的正实轴上，首个整数级共振恰发生于
\[
u=4,
\qquad
\lambda=1.
\]
\end{theorem}
\begin{proof}
命题 \ref{prop:real-input-40-output-spectral-collisions} 已给出
\[
Q\!\left(\frac{1}{\sqrt u},u\right)=\frac{2(\sqrt u-1)(\sqrt u+1)^2}{u^{3/2}},
\qquad
Q\!\left(-\frac{1}{\sqrt u},u\right)=\frac{2(\sqrt u-1)^2(\sqrt u+1)}{u^{3/2}},
\]
因而在 \(u>0\) 上，
\[
Q(\pm u^{-1/2},u)=0 \iff u=1.
\]
这说明在 \(1<u<4\) 内不存在原子分支与 \(Q\)-谱的碰撞。

同一命题又给出
\[
Q(1,u)=-u(u-4)(u-1)(u+1),
\qquad
Q(-1,u)=-(u-1)^2(u^2+6u-2).
\]
当 \(1<u<4\) 时，前式显然非零；而
\[
u^2+6u-2>0
\qquad (u>1),
\]
故后式亦非零。因此在 \(1<u<4\) 内不存在 \(Q\)-谱于 \(z=\pm1\) 的整数级共振。最后，\(u=4\) 时 \(Q(1,4)=0\)，从而给出首个正实整数级共振 \(\lambda=1\)。证毕。
\end{proof}

\begin{theorem}[主 Perron 根相对于原子谱支存在严格 \texorpdfstring{$\sqrt u$}{sqrt(u)} 级分离，并从上方逼近零温四态 Perron 常数]\label{thm:xi-time-part73-perron-vs-atomic-sqrtu-separation}
沿用推论 \ref{cor:real-input-40-factor-sqrtu-eigs}、命题 \ref{prop:real-input-40-pressure-freezing} 与命题 \ref{prop:real-input-40-output-mirror-branch} 的记号。记 \(\lambda(u)\) 为输出势的 Perron 正根，并令
\[
P(\theta):=\log\lambda(e^\theta),
\qquad
c^2=\rho(B_\infty),
\]
其中 \(c\approx 1.89635299351377\)、\(d\approx 0.807943317323508\)，且 \(c^2\) 是四次方程
\[
y^4-6y^3+9y^2-y-1=0
\]
的最大正根。则当 \(u\to+\infty\) 时有：
\[
\boxed{\;
\lambda(u)-\sqrt u
=
(c-1)\sqrt u+d+O(u^{-1/2}),\;}
\]
因而
\[
\boxed{\;
\lambda(u)-\sqrt u\longrightarrow +\infty.\;}
\]
进一步，
\[
\boxed{\;
\frac{\lambda(u)^2}{u}
=
\rho(B_\infty)+\frac{2cd}{\sqrt u}+O(u^{-1}),\;}
\]
特别地，由于 \(d>0\)，该比值从上方逼近 \(\rho(B_\infty)\)。
最后，若取 \(u=e^\theta\)，则
\[
\boxed{\;
P(\theta)-\frac{\theta}{2}
=
\log c+\frac{d}{c}e^{-\theta/2}+O(e^{-\theta})
\longrightarrow
\log c>0.\;}
\]
因此显式原子谱支 \(\pm\sqrt u\) 只决定零温端点的斜率 \(1/2\)，而主导自由能的幅度常数则完全由零温四态核心 \(B_\infty\) 决定。
\end{theorem}
\begin{proof}
由命题 \ref{prop:real-input-40-output-mirror-branch} 的 Puiseux 展开，
\[
\lambda(u)=c\sqrt u+d+\frac{e}{\sqrt u}+\frac{f}{u}+O(u^{-3/2}).
\]
与显式原子谱支 \(\sqrt u\) 相减，便得到
\[
\lambda(u)-\sqrt u
=
(c-1)\sqrt u+d+O(u^{-1/2}).
\]
由于 \(c>1\)，主项 \((c-1)\sqrt u\) 严格为正并趋于无穷，从而得到第一条。

再将同一展开平方：
\[
\lambda(u)^2
=
c^2u+2cd\sqrt u+O(1).
\]
两边同除以 \(u\)，并用 \(c^2=\rho(B_\infty)\)，即得
\[
\frac{\lambda(u)^2}{u}
=
\rho(B_\infty)+\frac{2cd}{\sqrt u}+O(u^{-1}).
\]
由于 \(d>0\)，该展开表明收敛从上方发生。

最后令 \(u=e^\theta\)。由
\[
\lambda(e^\theta)
=
c\,e^{\theta/2}
\left(
1+\frac{d}{c}e^{-\theta/2}+O(e^{-\theta})
\right)
\]
可得
\[
P(\theta)
=
\log\lambda(e^\theta)
=
\frac{\theta}{2}+\log c+\frac{d}{c}e^{-\theta/2}+O(e^{-\theta}),
\]
这就给出最后一条公式。证毕。
\end{proof}

\paragraph{二维切片塌缩与原子 Witt 校正的变量分裂}
在二维加权 \(\zeta(z;u,v)\) 的 \((2+8)\) 次显式分解、两条 rigid branches \(\pm\sqrt v\) 与一参数输出位势中的原子 Witt 剥离都已建立之后，还可把输出完全抑制切片与二维/一维原子分量的变量分工进一步压缩为下列两条接口。

\begin{theorem}[二维加权 \texorpdfstring{$\zeta(z;u,v)$}{zeta(z;u,v)} 在 \texorpdfstring{$v=0$}{v=0} 切片上的完全谱塌缩]\label{thm:xi-time-part73a-twoparam-vzero-total-spectral-collapse}
\leanverified{paper\_xi\_time\_part73a\_twoparam\_vzero\_total\_spectral\_collapse}
沿用定理 \ref{thm:real-input-40-det-uv-2plus8} 的记号。则对任意碰撞参数 \(u\) 都有严格恒等式
\[
\boxed{\;
\Delta(z;u,0)=1-z.\;}
\]
因此：
\begin{enumerate}
  \item \(M_{u,0}\) 的非零谱恰为单点 \(\{1\}\)；
  \item 由于 \(M_{u,0}\in M_{40}(\CC)\)，其零特征值的代数重数恰为 \(39\)；
  \item 相应加权 Artin--Mazur \(\zeta\) 函数满足
  \[
  \boxed{\;
  \zeta(z;u,0)=\frac{1}{1-z}.\;}
  \]
  \item 若定义 ghost 坐标与 primitive 生成函数
  \[
  a_n(u,0):=\Tr(M_{u,0}^n),
  \qquad
  P(z;u,0):=\sum_{n\ge 1}p_n(u,0)z^n,
  \]
  则有
  \[
  \boxed{\;
  a_n(u,0)=1\qquad(n\ge 1),\;}
  \]
  以及
  \[
  \boxed{\;
  P(z;u,0)=z,\qquad
  p_1(u,0)=1,\quad p_n(u,0)=0\ (n\ge 2).\;}
  \]
\end{enumerate}
换言之，一旦完全抑制输出位 \(e=1\)，二维模型中的碰撞权 \(u\) 在谱层面完全失明。
\end{theorem}
\begin{proof}
由定理 \ref{thm:real-input-40-det-uv-2plus8}，
\[
\Delta(z;u,v)=(v z^2-1)P_8(z;u,v),
\]
其中
\[
P_8(z;u,v)
=
(u^2 v^4-u^2 v^3)z^8+\cdots +(uv+5v)z^2+z-1.
\]
将 \(v=0\) 代入，则除末两项外其余各项全都消失，从而
\[
P_8(z;u,0)=z-1.
\]
于是
\[
\Delta(z;u,0)=(-1)(z-1)=1-z.
\]
这就给出第一条恒等式。

又由于该加权矩阵作用在真实输入 \(40\) 状态核上，故 \(M_{u,0}\) 为 \(40\times 40\) 矩阵。特征多项式的非平凡部分既然仅为 \(1-z\)，则唯一非零特征值只能是 \(1\)，其余 \(39\) 个特征值全为 \(0\)。于是
\[
\zeta(z;u,0)=\Delta(z;u,0)^{-1}=(1-z)^{-1}.
\]
从而
\[
\log\zeta(z;u,0)=\sum_{n\ge 1}\frac{z^n}{n},
\]
故 ghost 系数满足 \(a_n(u,0)=1\) 对所有 \(n\ge 1\) 成立；同一式也表明 primitive 生成函数恰为
\[
P(z;u,0)=z.
\]
于是 \(p_1(u,0)=1\) 且 \(p_n(u,0)=0\)（\(n\ge 2\)）。证毕。
\end{proof}

\begin{corollary}[二维/一维原子 Witt 校正的变量分裂与单层激活]\label{cor:xi-time-part73a-atomic-witt-bivariate-univariate-variable-rigidity}
分别考虑二维加权模型 \(\zeta(z;u,v)\) 与一参数输出位势模型 \(\zeta(z,u)\)。
\begin{enumerate}
  \item 在二维模型中，若定义 primitive 生成函数
  \[
  P(z;u,v):=\sum_{n\ge 1}p_n(u,v)z^n,
  \qquad
  P_{\mathrm{core}}(z;u,v):=\sum_{n\ge 1}p_n^{\mathrm{core}}(u,v)z^n,
  \]
  以及 ghost 坐标
  \[
  \log\zeta(z;u,v)=\sum_{n\ge 1}\frac{a_n(u,v)}{n}z^n,
  \qquad
  \log\zeta_{\mathrm{core}}(z;u,v)=\sum_{n\ge 1}\frac{a_n^{\mathrm{core}}(u,v)}{n}z^n,
  \]
  则有
  \[
  \boxed{\;
  P(z;u,v)=v z^2+P_{\mathrm{core}}(z;u,v),\;}
  \]
  \[
  \boxed{\;
  a_n(u,v)=
  \begin{cases}
  2v^{n/2}+a_n^{\mathrm{core}}(u,v), & 2\mid n,\\[4pt]
  a_n^{\mathrm{core}}(u,v), & 2\nmid n.
  \end{cases}\;}
  \]
  因而对任意 \(r\ge 1\) 与任意 \(n\ge 1\)，都有
  \[
  \partial_u^r\!\bigl(p_n(u,v)-p_n^{\mathrm{core}}(u,v)\bigr)\equiv 0,
  \qquad
  \partial_u^r\!\bigl(a_n(u,v)-a_n^{\mathrm{core}}(u,v)\bigr)\equiv 0.
  \]
  \item 在一参数输出位势模型中，若记主极点 \(z_\star(u)=\lambda(u)^{-1}\)，则原子 Witt 校正恰在 primitive 的 \(n=2\) 层激活，并满足
  \[
  \boxed{\;
  p_n(u)=p_n^{\mathrm{core}}(u)+u\,\mathbf 1_{\{n=2\}},\;}
  \]
  \[
  \boxed{\;
  a_n(u)=a_n^{\mathrm{core}}(u)+2u^{n/2}\mathbf 1_{\{2\mid n\}},\;}
  \]
  \[
  \boxed{\;
  \log\mathfrak M(u)-\log\mathfrak M_{\mathrm{core}}(u)=u\,z_\star(u)^2.\;}
  \]
  特别地，在无权点 \(u=1\) 上，
  \[
  \log\mathfrak M(1)-\log\mathfrak M_{\mathrm{core}}(1)
  =
  z_\star(1)^2
  =
  \varphi^{-4}.
  \]
\end{enumerate}
\end{corollary}
\begin{proof}
对二维模型，固定任意 \(u\) 并把 \(v\) 视为加权变量。由定理 \ref{thm:real-input-40-det-uv-2plus8}，
\[
\zeta(z;u,v)=\frac{1}{1-vz^2}\,\zeta_{\mathrm{core}}(z;u,v).
\]
于是定理 \ref{thm:xi-finite-euler-primitive-surgery-additivity} 在单原子情形 \((m,E,\ell)=(1,1,2)\) 下直接给出
\[
P(z;u,v)=v z^2+P_{\mathrm{core}}(z;u,v).
\]
另一方面，定理 \ref{thm:xi-time-part73-twoparam-atomic-ghost-u-blindness} 已证明
\[
a_n(u,v)-a_n^{\mathrm{core}}(u,v)=2v^{n/2}\mathbf 1_{\{2\mid n\}}.
\]
由于上述二维原子修正只含 \(v\) 而不含 \(u\)，故关于 \(u\) 的任意高阶导数恒为零。

对一参数模型，全部公式都已由定理 \ref{thm:xi-real-input40-evenlength-abel-shift-rigidity} 给出；最后在 \(u=1\) 处再用该定理中的专门化
\[
z_\star(1)=\varphi^{-2}
\]
便得到
\[
z_\star(1)^2=\varphi^{-4}.
\]
证毕。
\end{proof}

\endinput


\paragraph{固定参数 Möbius 反演下的偶长 necklace 校正}
在上文采用参数提升型 primitive 生成函数之外，若固定 \(u,v\) 并仅对长度指标施行 Möbius 反演，则二维原子因子的偶长校正还可进一步压缩为一条显式 necklace 闭式。

\begin{theorem}[二维加权 periodic 谱中的普适偶长原子校正]\label{thm:xi-time-part73c-periodic-evenlength-atomic-correction}
沿用定理 \ref{thm:real-input-40-det-uv-2plus8} 与定理 \ref{thm:xi-time-part73-nonrigid-trace-newton-cayley-closure} 的记号。记
\[
a_n(u,v):=\operatorname{Tr}(M_{u,v}^n),
\qquad
\mathcal Q_8(z;u,v):=-P_8(z;u,v),
\]
并由
\[
-\log \mathcal Q_8(z;u,v)=\sum_{n\ge 1}\frac{\widetilde a_n(u,v)}{n}z^n
\]
定义 residual periodic 系数 \(\widetilde a_n(u,v)\)。则对任意 \(n\ge 1\)，有
\[
a_n(u,v)=\widetilde a_n(u,v)+2v^{n/2}\mathbf 1_{\{2\mid n\}}.
\]
等价地，对每个 \(m\ge 0\) 有
\[
a_{2m+1}(u,v)=\widetilde a_{2m+1}(u,v),
\]
而对每个 \(m\ge 1\) 有
\[
a_{2m}(u,v)=\widetilde a_{2m}(u,v)+2v^m.
\]
特别地，该偶长校正与碰撞参数 \(u\) 完全无关。
\leanverified{evenLengthCorrection}
\leanverified{evenLengthCorrection\_odd}
\leanverified{evenLengthCorrection\_even}
\leanverified{evenLengthCorrection\_cases}
\leanverified{evenLengthCorrection\_eq\_zero\_iff}
\leanverified{evenLengthCorrection\_pos\_iff}
\leanverified{evenLengthCorrection\_zero}
\leanverified{evenLengthCorrection\_two}
\leanverified{paper\_zeta\_evenLength\_small\_values\_table}
\end{theorem}
\begin{proof}
由定理 \ref{thm:real-input-40-det-uv-2plus8}，
\[
\Delta(z;u,v)=(vz^2-1)P_8(z;u,v)=(1-vz^2)\mathcal Q_8(z;u,v).
\]
因此
\[
\zeta(z;u,v)=\frac{1}{\Delta(z;u,v)}
=
\frac{1}{1-vz^2}\,\mathcal Q_8(z;u,v)^{-1}.
\]
又由 \(\mathcal Q_8(0;u,v)=1\)，取形式对数得到
\[
\log\zeta(z;u,v)
=
-\log(1-vz^2)-\log \mathcal Q_8(z;u,v).
\]
另一方面，
\[
\log\zeta(z;u,v)=\sum_{n\ge 1}\frac{a_n(u,v)}{n}z^n,
\]
且
\[
-\log(1-vz^2)=\sum_{m\ge 1}\frac{v^m}{m}z^{2m}.
\]
把这三式逐项比较，便得
\[
\frac{a_n(u,v)-\widetilde a_n(u,v)}{n}
=
\begin{cases}
\dfrac{v^{n/2}}{n/2},& 2\mid n,\\[6pt]
0,& 2\nmid n,
\end{cases}
\]
亦即
\[
a_n(u,v)-\widetilde a_n(u,v)=2v^{n/2}\mathbf 1_{\{2\mid n\}}.
\]
最后一条只是同一公式的直接重述。证毕。
\end{proof}

\begin{corollary}[固定参数 Möbius 本原系数的偶长 necklace 校正]\label{cor:xi-time-part73c-fixed-parameter-necklace-correction}
固定参数 \(u,v\)。定义固定参数 Möbius 本原系数 \(p_n^{\mathrm{fix}}(u,v)\)、\(\widetilde p_n^{\mathrm{fix}}(u,v)\) 为
\[
a_n(u,v)=\sum_{d\mid n} d\,p_d^{\mathrm{fix}}(u,v),
\qquad
\widetilde a_n(u,v)=\sum_{d\mid n} d\,\widetilde p_d^{\mathrm{fix}}(u,v).
\]
则对任意 \(n\ge 1\)，有
\[
p_n^{\mathrm{fix}}(u,v)-\widetilde p_n^{\mathrm{fix}}(u,v)
=
\frac{2}{n}
\sum_{\substack{d\mid n\\ 2\mid n/d}}
\mu(d)\,v^{n/(2d)}.
\]
特别地：
\[
n\ \text{为奇数}
\quad\Longrightarrow\quad
p_n^{\mathrm{fix}}(u,v)=\widetilde p_n^{\mathrm{fix}}(u,v),
\]
\[
p_{2m}^{\mathrm{fix}}(u,v)-\widetilde p_{2m}^{\mathrm{fix}}(u,v)
=
\frac{1}{m}\sum_{d\mid m}\mu(d)\,v^{m/d}
\qquad (m\ge 1).
\]
若记
\[
M_m(v):=\frac{1}{m}\sum_{d\mid m}\mu(d)\,v^{m/d},
\]
则 \(M_m(v)\) 正是第 \(m\) 个 necklace 多项式。因而固定参数的 primitive 校正只支撑于偶长层，并在长度 \(2m\) 处恰等于 \(M_m(v)\)。
\leanverified{necklaceCorrectionKernel}
\leanverified{necklaceCorrectionKernel\_odd}
\leanverified{necklaceCorrectionKernel\_even}
\leanverified{necklaceCorrectionKernel\_odd\_eq\_zero}
\leanverified{necklaceCorrectionKernel\_divisor\_form}
\leanverified{necklaceCorrectionKernel\_values\_extended}
\leanverified{necklaceCorrectionKernel\_ternary}
\leanverified{paper\_necklaceCorrectionKernel\_extended}
\leanverified{paper\_xi\_time\_part73c\_fixed\_parameter\_necklace\_correction}
\leanverified{paper\_necklace\_correction\_two\_values}
\leanverified{necklaceCorrectionKernel\_at\_two\_pos\_bounded}
\leanverified{paper\_zeta\_evenLength\_correction\_package}
\leanverified{evenLength\_sub\_necklace\_at\_prime}
\leanverified{evenLengthCorrection\_partial\_sum}
\end{corollary}
\begin{proof}
由经典 Möbius 反演，
\[
n\,p_n^{\mathrm{fix}}(u,v)=\sum_{d\mid n}\mu(d)\,a_{n/d}(u,v),
\]
\[
n\,\widetilde p_n^{\mathrm{fix}}(u,v)=\sum_{d\mid n}\mu(d)\,\widetilde a_{n/d}(u,v).
\]
两式相减，并代入定理 \ref{thm:xi-time-part73c-periodic-evenlength-atomic-correction}，得到
\[
n\Bigl(
p_n^{\mathrm{fix}}(u,v)-\widetilde p_n^{\mathrm{fix}}(u,v)
\Bigr)
=
2\sum_{\substack{d\mid n\\ 2\mid n/d}}
\mu(d)\,v^{n/(2d)}.
\]
除以 \(n\) 即得第一式。

若 \(n\) 为奇数，则条件 \(2\mid n/d\) 不可能成立，故校正为零。若 \(n=2m\)，则
\[
d\mid 2m,\ 2\mid \frac{2m}{d}
\quad\Longleftrightarrow\quad
d\mid m,
\]
于是上式化为
\[
p_{2m}^{\mathrm{fix}}(u,v)-\widetilde p_{2m}^{\mathrm{fix}}(u,v)
=
\frac{1}{m}\sum_{d\mid m}\mu(d)\,v^{m/d}.
\]
最后一条只是把右端识别为标准 necklace 多项式。证毕。
\leanverified{paper\_necklaceCorrectionKernel\_odd\_zero}
\leanverified{necklaceCorrectionKernel\_v2\_seeds}
\leanverified{necklaceCorrectionKernel\_v2\_necklace\_number}
\leanverified{paper\_necklace\_correction\_golden\_mean\_seeds}
\end{proof}

\endinput


\noindent
以上这些结果表明，真实输入 \(40\) 状态核在二维原子层、非刚性残余谱、\(u=1\) 的局部激活层与 \(u\to+\infty\) 的零温核心层之间存在严格的多尺度分裂：一方面，二维加权 \(\zeta\) 的单原子因子在 ghost 层只向偶长度注入 \(2v^{n/2}\)，并对碰撞权 \(u\) 完全失明；去除两条 rigid branches \(\pm\sqrt v\) 之后，其余全部非刚性迹又自动坍缩为由单一八次 core 因子控制的 Newton--Cayley 递推系统，并在 \(v=1\) 超曲面与 \((u,v)=(0,1)\) 尖点上继续发生可量化的谱秩塌缩；当进一步压到输出完全抑制的切片 \(v=0\) 时，整个二维谱又严格塌缩为唯一一个长度 \(1\) 的 primitive 轨道。另一方面，在碰撞零温缩放下，整个二维核心因子进一步退化为同一原子的单参数模型 \(1-vw\)；在一参数输出位势上，同一原子 Witt 机制只在 primitive 的 \(n=2\) 层激活，并把 Abel 有限部刚性平移为 \(u z_\star(u)^2\)；再者，激活分支以无二次项的横截方式穿过 \(\lambda=0\)，其三次归一化 defect 在 \(u=1\) 处锁定为常数 \(3\)，去支后的剩余实谱则在最近退化点所围成的审计窗口内恢复为无碰撞解析族；而正实轴上的区间 \(1<u<4\) 又把原子碰撞与整数级共振共同隔离在外。最后，真正主导大 \(u\) 幅度常数的并不是显式原子谱支 \(\pm\sqrt u\)，而是零温四态核心 \(B_\infty\) 的 Perron 常数。因而在读出口径上，前述几条对应于中心投影式的相位盲读出，而其后的非刚性迹系统、局部极点、激活分支与谱壁垒则共同构成相位敏感的纤维内读出。

\paragraph{Lee--Yang 振幅域、有限支持账本障碍与 Smith--椭圆寄存器极值}
在 Stokes--Lee--Yang 三次数域的共同域识别、自然扩张的 prime-register 外置、Smith 信息亏损谱的离散曲率恢复，以及周期化 Bernoulli 共振与整数轴对齐椭圆半群的素数寄存器结构已经建立之后，还可进一步压缩出下列九条直接后果。它们分别刻画 Lee--Yang 振幅常数 \(B\) 的六次全复扩张、无限前史的有限支持 prime-slice 账本不可能性、固定总估值下 Smith 信息亏损的 sharp 极值原理，以及周期化 Fourier 共振与整数轴对齐椭圆在结构层和实现层上的刚性裂解。

\begin{theorem}[Lee--Yang 振幅常数 \texorpdfstring{$B$}{B} 的六次不可约性]\label{thm:xi-time-part74-leyang-amplitude-sixth-irreducibility}
\leanverified{paper\_xi\_time\_part74\_leyang\_amplitude\_sixth\_irreducibility}
沿用定理 \ref{thm:xi-terminal-zm-stokes-amplitude-square-cubic-rigidity} 与定理 \ref{thm:xi-terminal-zm-stokes-leyang-shared-artin-representation} 的记号，记
\[
b:=B^2,
\qquad
K:=\QQ(b)=\QQ(u)=\QQ(\lambda_*).
\]
则
\[
[\QQ(B):\QQ]=6,
\]
并且 \(B\) 在 \(\QQ\) 上满足的 primitive 不可约整式恰为
\[
\boxed{\;
162X^6+1593X^4+1998X^2+1184.\;}
\]
\end{theorem}
\begin{proof}
由定理 \ref{thm:xi-terminal-zm-stokes-amplitude-square-cubic-rigidity}，
\[
162b^3+1593b^2+1998b+1184=0
\]
是 \(b\) 在 \(\QQ\) 上的不可约三次式。再由定理
\ref{thm:xi-terminal-zm-stokes-leyang-shared-artin-representation}，
\[
K=\QQ(b)=\QQ(u)=\QQ(\lambda_*),
\qquad
[K:\QQ]=3.
\]

若 \(B\in K\)，则 \(b=B^2\) 为 \(K\) 中平方元，于是其范数必为 \(\QQ\) 中平方：
\[
N_{K/\QQ}(b)=N_{K/\QQ}(B)^2\in \QQ_{\ge 0}.
\]
另一方面，由三次极小多项式的首末系数可得
\[
N_{K/\QQ}(b)=(-1)^3\frac{1184}{162}=-\frac{592}{81},
\]
这在 \(\QQ\) 中不可能是平方，矛盾。故 \(B\notin K\)，从而 \(X^2-b\) 在 \(K[X]\) 中不可约，并且
\[
[\QQ(B):K]=2.
\]
于是
\[
[\QQ(B):\QQ]=2[K:\QQ]=6.
\]

显然 \(B\) 满足
\[
162B^6+1593B^4+1998B^2+1184=0.
\]
由于该整式次数为 \(6=[\QQ(B):\QQ]\)，故它就是 \(B\) 的 primitive 不可约整式。证毕。
\end{proof}

\begin{corollary}[Lee--Yang 振幅域的签名与二次扩张结构]\label{cor:xi-time-part74-leyang-amplitude-field-signature}
令
\[
L:=\QQ(B),
\qquad
K:=\QQ(B^2)=\QQ(b).
\]
则 \(L/K\) 为二次扩张；域 \(K\) 的签名为 \((1,1)\)，而 \(L\) 的签名为 \((0,3)\)。特别地，\(L\) 是全复六次域，而其立方子域 \(K\) 恰有一个实嵌入。
\end{corollary}
\leanverified{paper\_xi\_time\_part74\_leyang\_amplitude\_field\_signature}
\begin{proof}
由定理 \ref{thm:xi-time-part74-leyang-amplitude-sixth-irreducibility}，\([L:K]=2\)。再由定理 \ref{thm:xi-terminal-zm-stokes-amplitude-square-cubic-rigidity}，\(K\) 由不可约三次式
\[
f(t)=162t^3+1593t^2+1998t+1184
\]
生成，而推论 \ref{cor:xi-terminal-zm-stokes-b2-quadratic-resolvent} 给出的判别式平方自由部分为 \(-111\)，故 \(\Disc(f)<0\)。因此 \(K\) 的签名为 \((1,1)\)。

又因为 \(f\) 的系数全为正，所以 \(f(t)\) 在正实轴上无根；而三次实系数多项式必有实根，故其唯一实根只能是负数。于是对 \(K\) 的唯一实嵌入 \(\sigma:K\hookrightarrow \RR\)，有
\[
\sigma(b)<0.
\]
因此方程 \(x^2=\sigma(b)\) 在 \(\RR\) 中无解，故 \(\sigma\) 的任何延拓都不给出 \(L\) 的实嵌入。其余两条嵌入本来就是复嵌入，延拓后仍成共轭复对。由于 \([L:\QQ]=6\)，遂得 \(L\) 的签名为 \((0,3)\)。证毕。
\end{proof}

\begin{theorem}[自然扩张的有限支持 prime-slice 外置不可能性]\label{thm:xi-time-part74-natural-extension-finite-support-primeslice-impossibility}
\leanverified{paper\_xi\_time\_part74\_natural\_extension\_finite\_support\_primeslice\_impossibility}
设 \(T:X\twoheadrightarrow X\) 为可数集合上的有限纤维满射，并记其自然扩张为
\[
X^{\leftarrow}
:=
\left\{
(x_0,x_{-1},x_{-2},\dots)\in X^{\NN}:\ T(x_{-j})=x_{-j+1}\ \ (j\ge 1)
\right\}.
\]
对每个 \(N\ge 1\)，再记 \(N\)-步前史空间
\[
X^{\leftarrow,N}
:=
\left\{
(x_0,x_{-1},\dots,x_{-N})\in X^{N+1}:\ T(x_{-j})=x_{-j+1}\ \ (1\le j\le N)
\right\}.
\]
假设对每个 \(N\ge 1\) 以及每一层 \(0\le j\le N-1\)，都存在某个纤维点使该层纤维大小至少为 \(2\)。在定理 \ref{thm:killo-layered-prime-slices-finite-depth} 的分层 prime-slice 语义下，不存在任何单射
\[
\Psi:X^{\leftarrow}\longrightarrow X\times \ZZ_{>0},
\qquad
\Psi(\mathbf x)=\bigl(x_0,H(\mathbf x)\bigr),
\]
其编码分解为
\[
H(\mathbf x)=\prod_{j\ge 0}h_j(\mathbf x),
\]
并满足：
\begin{enumerate}
  \item 每个 \(h_j(\mathbf x)\) 的全部素因子都限制在预先固定的 prime slice \(\mathcal P_j\)；
  \item 存在统一常数 \(M\)，使得对所有 \(\mathbf x\in X^{\leftarrow}\) 与所有 \(j\ge M\) 都有
  \[
  h_j(\mathbf x)=1.
  \]
\end{enumerate}
\end{theorem}
\begin{proof}
由定理 \ref{thm:killo-natural-extension-branch-register}，自然扩张 \((X^{\leftarrow},\sigma)\) 与右移插入的分支寄存器系统共轭。另一方面，定理 \ref{thm:xi-time-part63b-infinite-depth-finite-prime-register-impossibility} 已经证明：只要每一层都存在真实分支，则任何保持全历史可恢复性的精确分层 prime-register 外置都不可能只使用有限多个 prime slices。

若存在所述编码 \(\Psi\) 与统一常数 \(M\)，则对每个 \(\mathbf x\in X^{\leftarrow}\)，整数 \(H(\mathbf x)\) 的全部非平凡素因子都落在前 \(M\) 个 prime slices 中；也就是说，这一外置实际上只使用了有限多个 prime slices。于是它正与定理 \ref{thm:xi-time-part63b-infinite-depth-finite-prime-register-impossibility} 的必要性结论矛盾。故这种有限支持外置不存在。证毕。
\end{proof}

\begin{corollary}[\texorpdfstring{$N$}{N} 步前史的最优活动 prime-slice 复杂度恰等于深度]\label{cor:xi-time-part74-finite-past-primeslice-complexity-depth}
\leanverified{paper\_xi\_time\_part74\_finite\_past\_primeslice\_complexity\_depth}
在上述设定下，假设对每个 \(N\ge 1\) 的前史空间 \(X^{\leftarrow,N}\)，其每一层都存在真实分支。记
\[
\mathfrak c_N
\]
为所有单射分层 prime-slice 外置
\[
\Phi_N:X^{\leftarrow,N}\hookrightarrow X\times \ZZ_{>0},
\qquad
\Phi_N(x_0,\dots,x_{-N})=(x_0,H_N(x_0,\dots,x_{-N}))
\]
在最坏情形下所需的最少非平凡 prime slices 数目。则
\[
\boxed{\;
\mathfrak c_N=N.\;}
\]
\end{corollary}
\begin{proof}
在 \(X^{\leftarrow,N}\) 上，每一层 \(0,\dots,N-1\) 都是活跃分支层。故定理 \ref{thm:xi-layered-primeslice-complexity-active-depth} 直接给出下界
\[
\mathfrak c_N\ge N.
\]
另一方面，定理 \ref{thm:killo-layered-prime-slices-finite-depth} 构造了恰在每一层写入一枚有效素数的单射外置，从而
\[
\mathfrak c_N\le N.
\]
两式合并即得 \(\mathfrak c_N=N\)。证毕。
\end{proof}

\begin{theorem}[固定总估值下 Smith 信息亏损的 sharp 极值原理]\label{thm:xi-time-part74-smith-loss-fixed-total-valuation-extremal}
\leanverified{paper\_xi\_time\_part74\_smith\_loss\_fixed\_total\_valuation\_extremal}
设 \(A\in M_{m\times n}(\ZZ)\) 满足 \(\operatorname{rank}_{\QQ}(A)=m\)，其 Smith 不变量为
\[
d_1\mid d_2\mid \cdots \mid d_m.
\]
固定素数 \(p\) 与整数 \(k\ge 1\)，记
\[
e_i:=v_p(d_i),
\qquad
E:=\sum_{i=1}^{m}e_i,
\qquad
f_{p,k}(A):=\sum_{i=1}^{m}\min(k,e_i).
\]
把 \(E\) 作 Euclidean 分解
\[
E=mq+r,
\qquad
0\le r<m.
\]
若只固定总估值 \(E\) 而允许 \((e_1,\dots,e_m)\) 在全部 Smith 估值谱中变化，则
\[
\boxed{\;
\min f_{p,k}(A)=\min(k,E),\;}
\]
\[
\boxed{\;
\max f_{p,k}(A)=(m-r)\min(k,q)+r\min(k,q+1).\;}
\]
极小值由集中谱
\[
(0,\dots,0,E)
\]
达到；极大值由最均衡谱
\[
(\underbrace{q,\dots,q}_{m-r},\underbrace{q+1,\dots,q+1}_{r})
\]
达到。
\end{theorem}
\begin{proof}
记
\[
\psi_k(t):=\min(k,t)
\qquad (t\in \ZZ_{\ge 0}).
\]
则
\[
f_{p,k}(A)=\sum_{i=1}^{m}\psi_k(e_i).
\]
函数 \(\psi_k\) 的离散导数
\[
\delta_k(t):=\psi_k(t)-\psi_k(t-1)
\qquad (t\ge 1)
\]
满足 \(\delta_k(t)=1\) 当 \(t\le k\)，而 \(\delta_k(t)=0\) 当 \(t>k\)。因此 \(\delta_k\) 单调不增，\(\psi_k\) 是离散凹函数。

现取任意两个坐标 \(a\le b-2\)。把它们替换为 \(a+1,b-1\) 时，总和保持不变，而函数值变化为
\[
\bigl(\psi_k(a+1)-\psi_k(a)\bigr)-\bigl(\psi_k(b)-\psi_k(b-1)\bigr)
=
\delta_k(a+1)-\delta_k(b)\ge 0.
\]
因此，每一次“均衡化”操作都不会减小 \(\sum_i\psi_k(e_i)\)。反复执行直到任意两坐标之差至多为 \(1\)，便必然到达唯一的最均衡谱
\[
(\underbrace{q,\dots,q}_{m-r},\underbrace{q+1,\dots,q+1}_{r}),
\]
从而得到极大值
\[
(m-r)\psi_k(q)+r\psi_k(q+1)
=(m-r)\min(k,q)+r\min(k,q+1).
\]

反向地，每一次把质量从较小坐标转移到较大坐标的“集中化”操作都不会增大 \(\sum_i\psi_k(e_i)\)。不断集中后，最终得到极端谱
\[
(0,\dots,0,E),
\]
其对应值为
\[
\psi_k(E)=\min(k,E).
\]
于是这给出极小值。

最后，二者都可由对角 Smith 形直接实现：取矩阵
\[
\operatorname{diag}(p^{e_1},\dots,p^{e_m})
\]
即可得到相应估值谱；若需固定列数 \(n\ge m\)，只需在右侧附加 \(n-m\) 个零列。故上下界都是 sharp。证毕。
\end{proof}

\begin{corollary}[一般模数下 Smith 信息亏损的逐素数 sharp 极值]\label{cor:xi-time-part74-general-modulus-kernel-loss-extremal}
\leanverified{paper\_xi\_time\_part74\_general\_modulus\_kernel\_loss\_extremal}
沿用定理 \ref{thm:xi-time-part74-smith-loss-fixed-total-valuation-extremal} 的记号。对任意模数
\[
b=\prod_p p^{k_p},
\]
记
\[
E_p:=v_p(d_1\cdots d_m)=mq_p+r_p,
\qquad
0\le r_p<m.
\]
若只固定每个素数方向的总估值 \(E_p\) 而允许 Smith 谱在各素数方向重新分配，则诱导映射
\[
A_b:(\ZZ/b\ZZ)^n\to(\ZZ/b\ZZ)^m
\]
满足 sharp bounds
\[
\boxed{\;
(n-m)\log b+\sum_p \min(k_p,E_p)\log p
\le
\log \#\ker(A_b),\;}
\]
\[
\boxed{\;
\log \#\ker(A_b)
\le
(n-m)\log b+\sum_p\Bigl[(m-r_p)\min(k_p,q_p)+r_p\min(k_p,q_p+1)\Bigr]\log p.\;}
\]
并且上下界都可由适当的 Smith 谱精确达到。
\end{corollary}
\begin{proof}
由推论 \ref{cor:killo-kernel-size-general-modulus}，
\[
\log \#\ker(A_b)
=
(n-m)\log b+\sum_p\left(\sum_{i=1}^{m}\min(k_p,e_i(p))\right)\log p,
\]
其中 \(e_i(p):=v_p(d_i)\)。对每个固定素数 \(p\)，在总估值
\[
\sum_{i=1}^{m}e_i(p)=E_p
\]
已定的条件下，定理 \ref{thm:xi-time-part74-smith-loss-fixed-total-valuation-extremal} 给出
\[
\min \sum_i\min(k_p,e_i(p))=\min(k_p,E_p),
\]
以及
\[
\max \sum_i\min(k_p,e_i(p))
=(m-r_p)\min(k_p,q_p)+r_p\min(k_p,q_p+1).
\]
逐素数求和便得到所示上下界。

Sharpness 也逐素数独立实现：对每个 \(p\)，任选达到极小或极大的非降估值谱 \((e_1(p),\dots,e_m(p))\)。再定义
\[
d_i:=\prod_p p^{e_i(p)}.
\]
由于每个素数方向上都保持非降序，故仍有
\[
d_1\mid d_2\mid \cdots \mid d_m.
\]
于是这些 \(d_i\) 确实给出某个 Smith 链，并且在各素数方向同时达到相应极值，从而整体界是 sharp。证毕。
\end{proof}

\begin{theorem}[临界周期化测度在 Fibonacci 双频上的正极限与非 Rajchman 性]\label{thm:xi-time-part74-periodized-bernoulli-fibonacci-limit-nonrajchman}
令
\[
q:=\varphi^{-1},
\qquad
\nu:=\bigl(\RR\to \mathbb T_2=\RR/(2\ZZ)\bigr)_\#\mu_q^{(2)}.
\]
再记
\[
c_{\mathrm{Fib}}:=\lim_{n\to\infty}C_\varphi(F_n)>0
\]
为推论 \ref{cor:fold-resonance-ladder-fib-luc-limits} 给出的 Fibonacci 方向极限常数。则
\[
\bigl|\widehat \nu(F_m)\bigr|\longrightarrow c_{\mathrm{Fib}},
\qquad
\bigl|\widehat \nu(F_{m+1})\bigr|\longrightarrow c_{\mathrm{Fib}}.
\]
特别地，\(\nu\) 不是 Rajchman 测度；因而 \(\nu\) 不可能对 \(\mathbb T_2\) 上的 Haar 测度绝对连续。
\end{theorem}
\begin{proof}
由定理 \ref{thm:fold-pisot-bernoulli-convolution-representation}，
\[
\widehat \nu(u)=\widehat{\mu_q^{(2)}}(\pi u)
\qquad (u\in \ZZ),
\]
并且
\[
\bigl|\widehat{\mu_q^{(2)}}(\pi u)\bigr|=C_\varphi(u).
\]
故对任意 \(m\ge 1\)，有
\[
\bigl|\widehat \nu(F_m)\bigr|=C_\varphi(F_m),
\qquad
\bigl|\widehat \nu(F_{m+1})\bigr|=C_\varphi(F_{m+1}).
\]
再由推论 \ref{cor:fold-resonance-ladder-fib-luc-limits}，
\[
C_\varphi(F_m)\to c_{\mathrm{Fib}},
\qquad
C_\varphi(F_{m+1})\to c_{\mathrm{Fib}},
\]
其中 \(c_{\mathrm{Fib}}>0\)。于是 \(\widehat \nu(u)\) 沿 Fibonacci 双频子列不趋于 \(0\)，从而 \(\nu\) 不是 Rajchman 测度。

若 \(\nu\) 对 Haar 测度绝对连续，则在有限测度空间 \(\mathbb T_2\) 上它必具有某个 \(L^1\) 密度；Riemann--Lebesgue 引理遂推出 \(\widehat \nu(u)\to 0\)，与上式矛盾。故 \(\nu\) 不可能绝对连续。证毕。
\end{proof}

\begin{corollary}[周期化临界测度全部卷积幂的共振刚性]\label{cor:xi-time-part74-periodized-bernoulli-convolution-powers-nonrajchman}
\leanverified{paper\_xi\_time\_part74\_periodized\_bernoulli\_convolution\_powers\_nonrajchman}
在定理 \ref{thm:xi-time-part74-periodized-bernoulli-fibonacci-limit-nonrajchman} 的设定下，对任意整数 \(r\ge 1\)，卷积幂 \(\nu^{*r}\) 满足
\[
\bigl|\widehat{\nu^{*r}}(F_m)\bigr|
\longrightarrow
c_{\mathrm{Fib}}^r>0.
\]
特别地，\(\nu^{*r}\) 对所有 \(r\ge 1\) 都不是 Rajchman 测度，也都不可能对 Haar 测度绝对连续。
\end{corollary}
\begin{proof}
Fourier 变换把卷积化为乘法：
\[
\widehat{\nu^{*r}}(u)=\widehat \nu(u)^r.
\]
由定理 \ref{thm:xi-time-part74-periodized-bernoulli-fibonacci-limit-nonrajchman}，
\[
\bigl|\widehat \nu(F_m)\bigr|\to c_{\mathrm{Fib}}>0.
\]
于是
\[
\bigl|\widehat{\nu^{*r}}(F_m)\bigr|
=
\bigl|\widehat \nu(F_m)\bigr|^r
\longrightarrow
c_{\mathrm{Fib}}^r>0.
\]
因此 \(\nu^{*r}\) 同样不是 Rajchman 测度。若它绝对连续，则其 Fourier 系数应当趋于 \(0\)，再由 Riemann--Lebesgue 引理得到矛盾。证毕。
\end{proof}

\begin{theorem}[整数轴对齐椭圆半群在结构层与实现层的半圆维裂解]\label{thm:xi-time-part74-integer-ellipse-structure-implementation-split}
令
\[
\mathsf E_{\NN}
:=
\left\{
E_{a,b}:\ \frac{x^2}{a^2}+\frac{y^2}{b^2}=1,\ a,b\in\NN_{>0}
\right\},
\qquad
E_{a,b}\odot E_{c,d}:=E_{ac,bd}.
\]
则：
\begin{enumerate}
  \item 作为交换幺半群，
  \[
  (\mathsf E_{\NN},\odot)\cong (\NN^{(\mathcal P)}\times \NN^{(\mathcal P)},+)
  \cong (\NN_{>0}\times \NN_{>0},\times);
  \]
  \item 其实现层半圆维满足
  \[
  \cdim^{\mathrm{impl}}(\mathsf E_{\NN})=\frac12;
  \]
  \item 其结构层同态半圆维满足
  \[
  \cdim^{\mathrm{hom}}_{+}(\mathsf E_{\NN},\odot)=+\infty.
  \]
\end{enumerate}
亦即，整数轴对齐椭圆半群在可实现性意义下仍只是可数对象，而在严格同态线性化意义下却需要可数无穷个彼此独立的素数轴。
\end{theorem}
\begin{proof}
第一条正是定理 \ref{thm:killo-ellipse-diagonal-prime-register-equivalence} 的内容。因而 \(\mathsf E_{\NN}\) 与 \(\NN_{>0}\times \NN_{>0}\) 等势，特别是一个可数无限集。由命题 \ref{prop:killo-implementation-vs-structural-half-circle-dimension} 的定义口径，可数无限集的实现层半圆维恰为 \(\frac12\)，从而得到第二条。

再看结构层。子幺半群
\[
\mathsf E_x:=\{E_{a,1}:a\in\NN_{>0}\}\subset \mathsf E_{\NN}
\]
满足
\[
(\mathsf E_x,\odot)\cong (\NN_{\times},\times).
\]
若存在某个有限 \(k\) 与单射幺半群同态
\[
\Phi:(\mathsf E_{\NN},\odot)\hookrightarrow (\NN^k,+),
\]
则其限制将给出
\[
\NN_{\times}\hookrightarrow (\NN^k,+)
\]
的单射幺半群同态，这与定理 \ref{thm:killo-prime-freedom-non-finitizability} 矛盾。故不存在任何有限秩加法寄存器对 \(\mathsf E_{\NN}\) 给出忠实同态实现，也即
\[
\cdim^{\mathrm{hom}}_{+}(\mathsf E_{\NN},\odot)=+\infty.
\]
证毕。
\end{proof}

\noindent
这九条结果把 Lee--Yang 分支振幅、自然扩张账本、Smith 信息亏损与椭圆 prime-register 几条后段链条进一步压缩成统一结构：一方面，Stokes 幅值平方虽与 Lee--Yang 三次共享同一立方子域，但其平方根振幅 \(B\) 必然逃离该三次数域，转而落入一个全复六次扩张；另一方面，有限深度前史虽然能够以逐层单素数的最优方式压缩，但这种最优性一旦进入无限前史便不再可有限支持化，从而把 prime-slice 账本刚性提升为真正的逆极限障碍；再者，Smith 信息亏损不仅由离散曲率完全恢复，而且在固定总估值约束下满足精确的集中--均衡极值律，并逐素数推送到一般模数的 sharp 核增长界；最后，Fibonacci 方向的周期化共振既在全部卷积幂上保持非 Rajchman 刚性，又与整数轴对齐椭圆半群的可数无穷素数轴结构一道，展示出实现层可数而结构层不可有限线性化的根本裂解。

\input{subsubsec__xi-time-protocol-conclusions-part75-hologram-shadow-productlaw-resonance-tail}

\endinput


\endinput


\endinput


\endinput
